% Generated mechanically from frozen input inputs/p03/ja/tex/dpa.tex.
% Do not edit this generated file; edit the bound locale source instead.

% OK, start here.
%

\setcounter{chapter}{22}
\chapter{分割冪代数}



\phantomsection
\label{dpa-section-phantom}


\section{序論}
\label{dpa-section-introduction}

\noindent
本章では分割冪代数とその応用について論じる。
参考文献として \cite{Berthelot} がある。





\section{分割冪}
\label{dpa-section-divided-powers}

\noindent
本節では分割冪環に関するいくつかの結果をまとめる。
空積は $1$ であるべきなので、規約 $0! = 1$ を用いる。

\begin{definition}
\label{dpa-definition-divided-powers}
$A$ を環、$I$ を $A$ のイデアルとする。写像の族
$\gamma_n : I \to I$（$n > 0$）が、すべての $n \geq 0$、$m > 0$、
$x, y \in I$ および $a \in A$ に対して
\begin{enumerate}
\item $\gamma_1(x) = x$、また $\gamma_0(x) = 1$ とおく。
\item $\gamma_n(x)\gamma_m(x) = \frac{(n + m)!}{n! m!} \gamma_{n + m}(x)$,
\item $\gamma_n(ax) = a^n \gamma_n(x)$,
\item $\gamma_n(x + y) = \sum_{i = 0, \ldots, n} \gamma_i(x)\gamma_{n - i}(y)$,
\item $\gamma_n(\gamma_m(x)) = \frac{(nm)!}{n! (m!)^n} \gamma_{nm}(x)$.
\end{enumerate}
を満たすとき、これを $I$ 上の{\it 分割冪構造}という。
\end{definition}

\noindent
定義に現れる有理数 $\frac{(n + m)!}{n! m!}$ および
$\frac{(nm)!}{n! (m!)^n}$ は実際には整数であることに注意する。
前者は $n + m$ 個の対象から $n$ 個を選ぶ方法の数であり、
後者は $nm$ 個の対象からなる集合を、それぞれ $m$ 個の対象をもつ
$n$ 個の組に分ける方法の数である。
以下では、この定義に関するいくつかの注意を述べ、$\gamma_n(x)$ が
$I$ における $x^n/n!$ の代用であることを示す。

\begin{lemma}
\label{dpa-lemma-silly}
$A$ を環、$I$ を $A$ のイデアルとする。
\begin{enumerate}
\item $\gamma$ が $I$ 上の分割冪構造であるならば\footnote{ここおよび以下では、
$\gamma$ は $I$ から $I$ への写像の列
$\gamma_1, \gamma_2, \gamma_3, \ldots$ の略記である。}、
$n \geq 1$、$x \in I$ に対して $n! \gamma_n(x) = x^n$ が成り立つ。
\end{enumerate}
$A$ は $\mathbf{Z}$-加群として無捩れであると仮定する。
\begin{enumerate}
\item[(2)] $I$ 上の分割冪構造は、存在すれば一意である。
\item[(3)] $\gamma_n : I \to I$ が写像であるとき、
$$
\gamma\text{ は分割冪構造である}
\Leftrightarrow
n! \gamma_n(x) = x^n\ \forall x \in I, n \geq 1.
$$
\item[(4)] イデアル $I$ が分割冪構造をもつための必要十分条件は、
イデアルとしての $I$ の生成元の集合 $x_i$ であって、すべての
$n \geq 1$ に対して $x_i^n \in (n!)I$ を満たすものが存在することである。
\end{enumerate}
\end{lemma}

\begin{proof}
\emph{(1) の証明。}
$\gamma$ が分割冪構造であるとする。条件 (2) において
$n,m$ の代わりに $1,n-1$ を用いると、
$n \gamma_n(x) = \gamma_1(x)\gamma_{n - 1}(x)$。
したがって、帰納法と条件 (1) により $n! \gamma_n(x) = x^n$ を得る。

\medskip\noindent
$A$ は $\mathbf{Z}$-加群として無捩れであると仮定する。
\emph{(2) の証明。} これは (1) から明らかである。

\medskip\noindent
\emph{(3) の証明。}
すべての $x \in I$ と $n \geq 1$ に対して
$n! \gamma_n(x) = x^n$ が成り立つと仮定する。
$A \subset A \otimes_{\mathbf{Z}} \mathbf{Q}$ であるから、定義
\ref{dpa-definition-divided-powers} の公理 (1)--(5) は、$A$ が
$\mathbf{Q}$-代数である場合に証明すれば十分である。
この場合 $\gamma_n(x) = x^n/n!$ であり、(1)--(5) は直接確認できる。
たとえば (4) は二項公式
$$
(x + y)^n = \sum_{i = 0, \ldots, n} \frac{n!}{i!(n - i)!} x^iy^{n - i}
$$
に対応する。係数が正しいことを確かめるため、読者にはこれらの確認を勧める。

\medskip\noindent
\emph{(4) の証明。}
$x_i$ がイデアルとしての $I$ の生成元であり、すべての $n \geq 1$
に対して $x_i^n \in (n!)I$ が成り立つと仮定する。すべての
$x \in I$ に対して $x^n \in (n!)I$ が成り立つと主張する。
この主張が成り立てば $\gamma_n(x) = x^n/n!$ とおくことができ、
(3) により、これは分割冪構造を与える。
主張を証明するため、まず $x = ax_i$ に対して成り立つことに注意する。
したがって、主張はアーベル群としての $I$ の生成元の集合について成り立つ。
これらの生成元による表示の長さに関する帰納法により、主張が $x$ と $y$
について成り立つときに $x+y$ についても成り立つことを示せば十分である。
これは二項定理から直ちに従う。
\end{proof}

\begin{example}
\label{dpa-example-ideal-generated-by-p}
$p$ を素数とする。$p$ で割り切れない任意の整数 $n$ が $A$ で
可逆であるような環 $A$、すなわち $\mathbf{Z}_{(p)}$-代数 $A$ をとる。
このとき $I = pA$ は標準的な分割冪構造をもつ。実際、
$x = pa \in I$ に対して
$$
\gamma_n(x) = \frac{p^n}{n!} a^n
$$
とおく。$n \geq 1$ に対して
$p^n/n! \in p\mathbf{Z}_{(p)}$ であることは直ちに確認できる
（たとえば、これは $n!$ の素因数分解における $p$ の指数が
$\left\lfloor n/p \right\rfloor + \left\lfloor n/p^2 \right\rfloor
+ \left\lfloor n/p^3 \right\rfloor + \ldots$ であることから従う）。
したがって、この定義は意味をもち、写像の列 $\gamma_n : I \to I$ を与える。
定義 \ref{dpa-definition-divided-powers} の条件 (1)--(5) が満たされることは
直接確認できる。別の見方として、この定義が
$A_0 = \mathbf{Z}_{(p)}$ に対して成り立つことは明らかであり、そこから
補題 \ref{dpa-lemma-gamma-extends} により結論が従う。
\end{example}

\noindent
$A$ の任意のイデアル $I$ と $I$ 上の任意の分割冪構造 $\gamma$ に対して、
$\gamma_n\left(0\right) = 0$ であることに注意する（これは定義
\ref{dpa-definition-divided-powers} の公理 (3) に $a=0$ を適用すれば従う）。

\begin{lemma}
\label{dpa-lemma-check-on-generators}
$A$ を環、$I$ を $A$ のイデアルとし、
$\gamma_n : I \to I$（$n \geq 1$）を写像の列とする。次を仮定する。
\begin{enumerate}
\item[(a)] 定義 \ref{dpa-definition-divided-powers} の (1)、(3)、(4) が
すべての $x,y\in I$ に対して成り立つ。
\item[(b)] 性質 (2) と (5) が、イデアルとしての $I$ のある生成元集合の
各 $x$ に対して成り立つ。
\end{enumerate}
このとき $\gamma$ は $I$ 上の分割冪構造である。
\end{lemma}

\begin{proof}
本証明の番号 (1)、(2)、(3)、(4)、(5) は、定義
\ref{dpa-definition-divided-powers} に列挙された条件を指す。
(3) を適用すると、(2) と (5) が $x$ に対して成り立てば、すべての
$a\in A$ に対して $ax$ についても成り立つことが分かる。
したがって (b) より、(2) と (5) はアーベル群としての $I$ の
ある生成元集合について成り立つ。これらの生成元による表示の長さに関する
帰納法により、(2) と (5) が $x,y\in I$ に対して成り立つとき、
$x+y$ に対しても成り立つことを示せば十分である。

\medskip\noindent
\emph{$x+y$ に対する性質 (2)。} (4) により、
$$
\gamma_n(x + y)\gamma_m(x + y) =
\sum\nolimits_{i + j = n,\ k + l = m}
\gamma_i(x)\gamma_k(x)\gamma_j(y)\gamma_l(y)
$$
$x$ と $y$ に対する性質 (2) を用いると、これは
$$
\sum \frac{(i + k)!}{i!k!}\frac{(j + l)!}{j!l!}
\gamma_{i + k}(x)\gamma_{j + l}(y)
$$
に等しい。これを展開式
$$
\gamma_{n + m}(x + y) = \sum \gamma_a(x)\gamma_b(y)
$$
と比較すると、$a+b=n+m$ が与えられたとき
$$
\sum\nolimits_{i + k = a,\ j + l = b,\ i + j = n,\ k + l = m}
\frac{(i + k)!}{i!k!}\frac{(j + l)!}{j!l!}
=
\frac{(n + m)!}{n!m!}.
$$
を示せば十分である。これを直接論じる代わりに、多項式環
$\mathbf{Q}[x,y]$ のイデアル $I=(x,y)$ では、$f\in I$ に対する
$\gamma_n(f)=f^n/n!$ が $I$ 上の分割冪構造を定めるので、この等式が
成り立つことに注意する。したがって、上の有理数の等式は成り立つ。

\medskip\noindent
\emph{$x+y$ に対する性質 (5)。}
(1)--(4) が成り立ち、さらに (5) が $x$ と $y$ に対して成り立つと仮定する。
再び証明を有理数の等式に帰着する。(4) を用いると
$\gamma_n(\gamma_m(x + y)) = \gamma_n(\sum \gamma_i(x)\gamma_j(y))$.
と書ける。もう一度 (4) を用いると、$\gamma_n(\gamma_m(x+y))$ は、
各項が $\gamma_k(\gamma_i(x)\gamma_j(y))$ の形の因子の積であるような
項の和として書ける。$i>0$ ならば
\begin{align*}
\gamma_k(\gamma_i(x)\gamma_j(y)) & =
\gamma_j(y)^k\gamma_k(\gamma_i(x)) \\
& = \frac{(ki)!}{k!(i!)^k} \gamma_j(y)^k \gamma_{ki}(x) \\
& =
\frac{(ki)!}{k!(i!)^k} \frac{(kj)!}{(j!)^k} \gamma_{ki}(x) \gamma_{kj}(y)
\end{align*}
である。ここで第1の等号では (3)、第2の等号では $x$ に対する (5)、
第3の等号では (2) をちょうど $k$ 回用いた。$y$ に対する (5) を用いると、
同じ等式が $i=0$ の場合にも成り立つ。このように (5) 以外のすべての
公理を用いていくと、
$$
\gamma_n(\gamma_m(x + y)) =
\sum\nolimits_{i + j = nm} c_{ij}\gamma_i(x)\gamma_j(y)
$$
と書ける。ここで $c_{ij}\in\mathbf{Z}$ はある普遍定数である。
この等式が多項式環 $\mathbf{Q}[x,y]$ で成り立つことから、再び、
すべての係数 $c_{ij}$ が $(nm)!/n!(m!)^n$ に等しいことが従う。
\end{proof}

\begin{lemma}
\label{dpa-lemma-two-ideals}
$A$ を環とし、二つのイデアル $I,J\subset A$ をとる。
$\gamma$ を $I$ 上の分割冪構造、$\delta$ を $J$ 上の分割冪構造とする。
このとき
\begin{enumerate}
\item $\gamma$ と $\delta$ は $IJ$ 上で一致する。
\item $\gamma$ と $\delta$ が $I\cap J$ 上で一致するならば、これらは
$I+J$ 上の一意な分割冪構造 $\epsilon$ の制限である。
\end{enumerate}
\end{lemma}

\begin{proof}
$x\in I$、$y\in J$ とする。このとき
$$
\gamma_n(xy) = y^n\gamma_n(x) = n! \delta_n(y) \gamma_n(x) =
\delta_n(y) x^n = \delta_n(xy).
$$
である。したがって、$\gamma$ と $\delta$ は $IJ$ の（加法的）
生成元の集合上で一致する。定義 \ref{dpa-definition-divided-powers} の
性質 (4) より、これらは $IJ$ 全体で一致する。

\medskip\noindent
$\gamma$ と $\delta$ が $I\cap J$ 上で一致すると仮定する。
$z\in I+J$ をとり、$x\in I$、$y\in J$ により $z=x+y$ と書く。
すべての $n\geq 1$ に対して
$$
\epsilon_n(z) = \sum \gamma_i(x)\delta_{n - i}(y)
$$
とおく。これが良定義であることを示すため、$x'\in I$、$y'\in J$ による
別の表示 $z=x'+y'$ をとる。このとき
$w=x-x'=y'-y\in I\cap J$ である。したがって
\begin{align*}
\sum\nolimits_{i + j = n} \gamma_i(x)\delta_j(y)
& =
\sum\nolimits_{i + j = n} \gamma_i(x' + w)\delta_j(y) \\
& =
\sum\nolimits_{i' + l + j = n} \gamma_{i'}(x')\gamma_l(w)\delta_j(y) \\
& =
\sum\nolimits_{i' + l + j = n} \gamma_{i'}(x')\delta_l(w)\delta_j(y) \\
& =
\sum\nolimits_{i' + j' = n} \gamma_{i'}(x')\delta_{j'}(y + w) \\
& =
\sum\nolimits_{i' + j' = n} \gamma_{i'}(x')\delta_{j'}(y')
\end{align*}
となり、所望の等式を得る。これにより、すべての $n\geq 1$ に対して
写像 $\epsilon_n:I+J\to I+J$ を定義した。明らかに
$\epsilon_n\mid_I=\gamma_n$ かつ $\epsilon_n\mid_J=\delta_n$ である。
写像の族 $\epsilon_n$ が定義 \ref{dpa-definition-divided-powers} の
条件 (1)--(5) を満たすことを確認する。性質 (1) と (3) は明らかである。
(4) を確認するため、$x,x'\in I$、$y,y'\in J$ によって
$z=x+y$、$z'=x'+y'$ と書き、次のように計算する。
\begin{align*}
\epsilon_n(z + z') & =
\sum\nolimits_{a + b = n} \gamma_a(x + x')\delta_b(y + y') \\
& =
\sum\nolimits_{i + i' + j + j' = n}
\gamma_i(x) \gamma_{i'}(x')\delta_j(y)\delta_{j'}(y') \\
& =
\sum\nolimits_{k = 0, \ldots, n}
\sum\nolimits_{i+j=k} \gamma_i(x)\delta_j(y)
\sum\nolimits_{i'+j'=n-k} \gamma_{i'}(x')\delta_{j'}(y') \\
& =
\sum\nolimits_{k = 0, \ldots, n}\epsilon_k(z)\epsilon_{n-k}(z')
\end{align*}
これは所望の等式である。補題 \ref{dpa-lemma-check-on-generators} により、
あとは $I$ または $J$ の元について (2) と (5) を示せば十分である。
$\gamma$ と $\delta$ はいずれも分割冪構造であるから、これは明らかである。

\medskip\noindent
以上により、$I$ および $J$ への制限がそれぞれ $\gamma$ および
$\delta$ であるような $I+J$ 上の分割冪構造 $\epsilon$ の存在が示された。
一意性も明らかである。
\end{proof}

\begin{lemma}
\label{dpa-lemma-nil}
$p$ を素数、$A$ を環、$I\subset A$ をイデアルとし、
$\gamma$ を $I$ 上の分割冪構造とする。$p$ は $A/I$ において
冪零であると仮定する。このとき、$I$ が局所冪零であるための必要十分条件は、
$p$ が $A$ において冪零であることである。
\end{lemma}

\begin{proof}
$A$ において $p^N=0$ であるならば、$(pN)!$ は $p^N$ で割り切れるので、
$x\in I$ に対して
$x^{pN}=(pN)!\gamma_{pN}(x)=0$ である。
逆に、$I$ が局所冪零であると仮定する。さらに $p$ は $A/I$ において
冪零であると仮定しているので、ある $r$ に対して $p^r\in I$ である。
したがって $p^r$ は冪零であり、ゆえに $p$ も冪零である。
\end{proof}








\section{分割冪環}
\label{dpa-section-divided-power-rings}

\noindent
分割冪環の圏が存在する。その定義は次のとおりである。

\begin{definition}
\label{dpa-definition-divided-power-ring}
{\it 分割冪環}とは、$A$ が環、
$I\subset A$ がイデアル、$\gamma=(\gamma_n)_{n\geq 1}$ が
$I$ 上の分割冪構造であるような三つ組 $(A,I,\gamma)$ のことである。
{\it 分割冪環の準同型}
$\varphi:(A,I,\gamma)\to(B,J,\delta)$ とは、$\varphi(I)\subset J$ を満たし、
すべての $x\in I$ と $n\geq 1$ に対して
$\delta_n(\varphi(x))=\varphi(\gamma_n(x))$ が成り立つような環準同型
$\varphi:A\to B$ のことである。
\end{definition}

\noindent
ときに「$(B,J,\delta)$ を $(A,I,\gamma)$ 上の分割冪代数とする」という。
これは、$(B,J,\delta)$ が分割冪環であり、分割冪環の準同型
$(A,I,\gamma)\to(B,J,\delta)$ が備わっていることを意味する。

\begin{lemma}
\label{dpa-lemma-limits}
分割冪環の圏はすべての極限をもち、それらは環の圏における極限と一致する。
\end{lemma}

\begin{proof}
空極限は零環である（これは奇妙であるが、必要である）。
分割冪環の族 $(A_t,I_t,\gamma_t)$（$t\in T$）の積は
$(\prod A_t,\prod I_t,\gamma)$ で与えられる。ここで
$\gamma_n((x_t))=(\gamma_{t,n}(x_t))$ である。
二つの準同型 $\alpha,\beta:(A,I,\gamma)\to(B,J,\delta)$ の等化子は、
イデアル $C\cap I$ とその上に誘導される分割冪を備えた
$C=\{a\in A\mid\alpha(a)=\beta(a)\}$ にほかならない。
したがって、すべての極限が存在する。Categories, Lemma
\ref{categories-lemma-limits-products-equalizers} を参照せよ。
\end{proof}

\noindent
次の補題は、分割冪代数の場合に非常に一般的な圏論的現象を例示する。

\begin{lemma}
\label{dpa-lemma-a-version-of-brown}
$\mathcal{C}$ を分割冪環の圏とし、
$F:\mathcal{C}\to\textit{Sets}$ を関手とする。次を仮定する。
\begin{enumerate}
\item ある基数 $\kappa$ が存在して、任意の $f\in F(A,I,\gamma)$ に対し、
$|A'|\leq\kappa$ であるような $\mathcal{C}$ の射
$(A',I',\gamma')\to(A,I,\gamma)$ と元
$f'\in F(A',I',\gamma')$ が存在し、$f$ は $f'$ の像である。
\item $F$ は極限と可換する。
\end{enumerate}
このとき $F$ は表現可能である。すなわち、$\mathcal{C}$ の対象
$(B,J,\delta)$ が存在して、
$$
F(A, I, \gamma) = \Hom_\mathcal{C}((B, J, \delta), (A, I, \gamma))
$$
が $(A,I,\gamma)$ に関して関手的に成り立つ。
\end{lemma}

\begin{proof}
これは Categories, Lemma \ref{categories-lemma-a-version-of-brown}
の特別な場合である。
\end{proof}

\begin{lemma}
\label{dpa-lemma-colimits}
分割冪環の圏はすべての余極限をもつ。
\end{lemma}

\begin{proof}
空余極限は、分割冪イデアル $(0)$ を備えた $\mathbf{Z}$ である。
一般の余極限について論じる。$\mathcal{C}$ を圏とし、
$c\mapsto(A_c,I_c,\gamma_c)$ を図式とする。関手
$$
F(B, J, \delta) = \lim_{c \in \mathcal{C}}
Hom((A_c, I_c, \gamma_c), (B, J, \delta))
$$
を考える。任意の $f=(f_c)_{c\in C}\in F(B,J,\delta)$ は次の性質をもつ。
像 $f_c(A_c)$ のすべては、濃度がある基数 $\kappa$ で抑えられる
$B$ の部分環 $B'$ を生成し、像 $f_c(I_c)$ のすべては $B'$ の
分割冪部分イデアル $J'$ を生成する。したがって $f$ は、$F(B')$ の元
$f'$ と包含 $B'\to B$ の合成として分解する。さらに $F$ は極限と可換する。
ゆえに補題 \ref{dpa-lemma-a-version-of-brown} を適用でき、$F$ は表現可能である。
\end{proof}

\begin{remark}
\label{dpa-remark-pushout}
次を考える。
$$
\xymatrix{
(B, J, \delta) \ar[r] & (B'', J'', \delta'') \\
(A, I, \gamma) \ar[r] \ar[u] & (B', J', \delta') \ar[u]
}
$$
これが分割冪環の圏における押し出しであるとする（押し出しは補題
\ref{dpa-lemma-colimits} により存在する）。このとき
\begin{enumerate}
\item $B''/J'' = B/J \otimes_{A/I} B'/J'$,
\item 写像 $\varphi : B \otimes_A B' \to B''$ は全射である。
\item $\varphi$ によって誘導される写像
$J \otimes_A B' + B \otimes_A J' \to J''$ は全射である。
\end{enumerate}
(1) を示すには、写像
$(B'', J'', \delta'') \to (C, (0), \emptyset)$ を考え、押し出しの定義を用いる。
(2) を示すには、像 $\tilde B = \varphi(B \otimes_A B')$ とイデアル
$\tilde J = \varphi(J \otimes_A B' + B \otimes_A J')$ を考える。
$\delta''$ は $\delta$ および $\delta'$ と両立するので、すべての
$n \geq 1$ に対して $\delta''_n(\tilde J) \subset \tilde J$ であることは明らかである。
したがって、押し出しの普遍性により
$\tilde B = B''$ かつ $\tilde J = J''$ でなければならない。
\end{remark}

\begin{remark}
\label{dpa-remark-forgetful}
忘却関手 $(A, I, \gamma) \mapsto A$ は余極限と可換しない。たとえば、
$$
\xymatrix{
(B, J, \delta) \ar[r] & (B'', J'', \delta'') \\
(A, I, \gamma) \ar[r] \ar[u] & (B', J', \delta') \ar[u]
}
$$
を、注意 \ref{dpa-remark-pushout} で論じた分割冪環の圏における押し出しとする。
このとき一般に写像 $B \otimes_A B' \to B''$ は同型ではない。
具体例は次のとおりである。
$(A, I, \gamma) = (\mathbf{Z}, (0), \emptyset)$,
$(B, J, \delta) = (\mathbf{Z}/4\mathbf{Z}, 2\mathbf{Z}/4\mathbf{Z}, \delta)$,
および
$(B', J', \delta') =
(\mathbf{Z}/4\mathbf{Z}, 2\mathbf{Z}/4\mathbf{Z}, \delta')$
ここで $\delta_2(2) = 2$、$\delta'_2(2) = 0$ である。
より正確には、補題 \ref{dpa-lemma-need-only-gamma-p} を用いて、$\delta$ および
$\delta'$ を、それぞれ $J$ および $J'$ 上の一意な分割冪構造であって、
$\delta_2 : J \to J$ および $\delta'_2 : J' \to J'$ がそれぞれ写像
$0 \mapsto 0, 2 \mapsto 2$ および $0 \mapsto 0, 2 \mapsto 0$ であるものとする。
すると $(B'', J'', \delta'') = (\mathbf{F}_2, (0), \emptyset)$ となり、
これはテンソル積と一致しない。
\end{remark}


\section{分割冪構造の延長}
\label{dpa-section-extend}

\noindent
定義は次のとおりである。

\begin{definition}
\label{dpa-definition-extends}
分割冪環 $(A, I, \gamma)$ と環準同型 $A \to B$ が与えられたとする。
$IB$ 上の分割冪構造 $\bar \gamma$ であって、
$(A, I, \gamma) \to (B, IB, \bar\gamma)$ が分割冪環の準同型となるものが
存在するとき、$\gamma$ は $B$ に{\it 延長される}という。
\end{definition}

\begin{lemma}
\label{dpa-lemma-gamma-extends}
$(A, I, \gamma)$ を分割冪環、$A \to B$ を環準同型とする。
$\gamma$ が $B$ に延長されるならば、その延長は一意である。
次の条件のうち少なくとも一つが成り立つと仮定する。
\begin{enumerate}
\item $IB = 0$,
\item $I$ は単項イデアルである。
\item $A \to B$ は平坦である。
\end{enumerate}
このとき $\gamma$ は $B$ に延長される。
\end{lemma}

\begin{proof}
$IB$ の任意の元は、$b_i \in B$、$x_i \in I$ を用いて有限和
$\sum\nolimits_{i=1}^t b_ix_i$ と書ける。$\gamma$ が $IB$ 上の
$\bar\gamma$ に延長されるならば、
$\bar\gamma_n(x_i) = \gamma_n(x_i)$ である。
したがって、定義 \ref{dpa-definition-divided-powers} の条件 (3) と (4) より
$$
\bar\gamma_n(\sum\nolimits_{i=1}^t b_ix_i) =
\sum\nolimits_{n_1 + \ldots + n_t = n}
\prod\nolimits_{i = 1}^t b_i^{n_i}\gamma_{n_i}(x_i)
$$
を得る。よって $\bar\gamma$ は、存在すれば一意である。

\medskip\noindent
$IB = 0$ ならば $\bar\gamma_n(0) = 0$ とおけばよい。$I = (x)$ ならば、
$\bar\gamma_n(bx) = b^n\gamma_n(x)$ と定義する。これは良定義である。
実際、$b'x = bx$、すなわち $(b - b')x = 0$ ならば、
\begin{align*}
b^n\gamma_n(x) - (b')^n\gamma_n(x)
& =
(b^n - (b')^n)\gamma_n(x) \\
& =
(b^{n - 1} + \ldots + (b')^{n - 1})(b - b')\gamma_n(x) = 0
\end{align*}
となる。$\gamma_n(I) \subset I$ であるから $\gamma_n(x)$ は $x$ で
割り切れ、したがって $b-b'$ により零化されるからである。
次に、定義 \ref{dpa-definition-divided-powers} の条件 (1)--(5) を示す。
(1)、(2)、(3)、(5) は構成から明らかである。
(4) について、$y,z\in IB$ を $y=bx$、$z=cx$ と書く。このとき
$y+z=(b+c)x$ であるから、
\begin{align*}
\bar\gamma_n(y + z)
& =
(b + c)^n\gamma_n(x) \\
& =
\sum \frac{n!}{i!(n - i)!}b^ic^{n -i}\gamma_n(x) \\
& =
\sum b^ic^{n - i}\gamma_i(x)\gamma_{n - i}(x) \\
& =
\sum \bar\gamma_i(y)\bar\gamma_{n -i}(z)
\end{align*}
となり、所望の等式を得る。

\medskip\noindent
$A \to B$ が平坦であると仮定する。$b_1, \ldots, b_r \in B$ および
$x_1, \ldots, x_r \in I$ をとる。$\bar\gamma_n$ が存在するならば、
$$
\bar\gamma_n(\sum b_ix_i) =
\sum b_1^{e_1} \ldots b_r^{e_r} \gamma_{e_1}(x_1) \ldots \gamma_{e_r}(x_r)
$$
となる。ここで和は $e_1 + \ldots + e_r = n$ を走る。
次に、$b_i = \sum a_{ij}c_j$ を満たす
$c_1, \ldots, c_s \in B$ と $a_{ij} \in A$ があると仮定する。
$y_j = \sum a_{ij}x_i$ とおく。このとき
$$
\sum b_1^{e_1} \ldots b_r^{e_r} \gamma_{e_1}(x_1) \ldots \gamma_{e_r}(x_r) =
\sum c_1^{d_1} \ldots c_s^{d_s} \gamma_{d_1}(y_1) \ldots \gamma_{d_s}(y_s)
$$
が $B$ において成り立つと主張する。ここで右辺の和は
$d_1 + \ldots + d_s = n$ を走る。実際、分割冪構造の公理を用いると、
両辺を、係数が $\mathbf{Z}[a_{ij}]$ に属し、各項が
$c_1^{d_1} \ldots c_s^{d_s}\gamma_{e_1}(x_1) \ldots \gamma_{e_r}(x_r)$。
の形であるような項の和に展開できる。係数が一致することを見るには、
$\gamma$ を $(x_1, \ldots, x_r)$ 上の一意な分割冪構造としたとき、この結果が
$\mathbf{Q}[x_1, \ldots, x_r, c_1, \ldots, c_s, a_{ij}]$ で成り立つことに
注意すればよい。Lazard の定理（Algebra, Theorem
\ref{algebra-theorem-lazard}）により、$B$ を有限自由 $A$-加群の
有向余極限として書ける。特に、$z\in IB$ が
$z=\sum x_ib_i$ とも $z=\sum x'_{i'}b'_{i'}$ とも書けるならば、
$b_i=\sum a_{ij}c_j$、$b'_{i'}=\sum a'_{i'j}c_j$、および
$y_j = \sum x_ia_{ij} = \sum x'_{i'}a'_{i'j}$ を満たす
$c_1,\ldots,c_s\in B$ と $a_{ij},a'_{i'j}\in A$ を見いだせる\footnote{これは
Algebra, Theorem \ref{algebra-theorem-lazard} を用いずに証明することもできる。
実際、$z = \sum x_ib_i$ かつ $z = \sum x'_{i'}b'_{i'}$ ならば、
$\sum x_ib_i - \sum x'_{i'}b'_{i'} = 0$ は $A$-加群 $B$ における関係式である。
したがって、Algebra, Lemma \ref{algebra-lemma-flat-eq} を、$f_i$ の代わりに
$x_i$ と $x'_{i'}$、$x_i$ の役割に $b_i$ と $b'_{i'}$ を用いて適用すると、
所望の $c_1, \ldots, c_s \in B$ および
$a_{ij}, a'_{i'j} \in A$ の存在を得る。}。
したがって、上の手順は $IB$ 上の良定義な写像 $\bar\gamma_n$ を与える。
構成により $\bar\gamma$ は条件 (1)、(3)、(4) を満たす。さらに
$x\in I$ に対して $\bar\gamma_n(x)=\gamma_n(x)$ である。よって補題
\ref{dpa-lemma-check-on-generators} により、$\bar\gamma$ は $IB$ 上の分割冪構造である。
\end{proof}

\begin{lemma}
\label{dpa-lemma-kernel}
$(A, I, \gamma)$ を分割冪環とする。
\begin{enumerate}
\item $\varphi : (A, I, \gamma) \to (B, J, \delta)$ が分割冪環の準同型ならば、
すべての $n\geq 1$ に対して $\Ker(\varphi) \cap I$ は $\gamma_n$ で保たれる。
\item $\mathfrak a \subset A$ をイデアルとし、
$I' = I \cap \mathfrak a$ とおく。次の条件は同値である。
\begin{enumerate}
\item すべての $n > 0$ に対して $I'$ は $\gamma_n$ で保たれる。
\item $\gamma$ は $A/\mathfrak a$ に延長される。
\item イデアルとしての $I'$ の生成元の集合 $x_i$ であって、すべての
$n>0$ に対して $\gamma_n(x_i) \in I'$ を満たすものが存在する。
\end{enumerate}
\end{enumerate}
\end{lemma}

\begin{proof}
(1) の証明。これは明らかである。(2)(a) を仮定する。$x\in I$ に対して
$\bar\gamma_n(x \bmod I') = \gamma_n(x) \bmod I'$ と定義する。
これは良定義である。実際、定義 \ref{dpa-definition-divided-powers} (4) と、
仮定により $\gamma_j(y)\in I'$ であることから、$y\in I'$ に対して
$\gamma_n(x+y)=\gamma_n(x)\bmod I'$ となる。$\gamma$ が分割冪構造なので、
$\bar\gamma$ も明らかに分割冪構造である。したがって (2)(b) が成り立つ。
また (1) より (2)(b) は (2)(a) を含意する。(2)(a) が (2)(c) を
含意することは明らかである。(2)(c) を仮定する。
$x=ax_i$ ならば $\gamma_n(x)=a^n\gamma_n(x_i)\in I'$ であることに注意する。
したがって、アーベル群としての $I'$ のある生成元集合について
$\gamma_n(x)\in I'$ が成り立つ。これらの生成元による表示の長さに関する
帰納法により、$\forall n : \gamma_n(x), \gamma_n(y) \in I'$ ならば
$\forall n : \gamma_n(x+y)\in I'$ であることを示せば十分である。
これは分割冪構造の第4公理から直ちに従う。
\end{proof}

\begin{lemma}
\label{dpa-lemma-sub-dp-ideal}
$(A, I, \gamma)$ を分割冪環とし、$E\subset I$ を部分集合とする。
$\gamma$ で保たれ、かつすべての $f\in E$ を含む最小のイデアル
$J\subset I$ は、$\gamma_n(f)$（$n\geq 1$、$f\in E$）で生成される
イデアル $J$ である。
\end{lemma}

\begin{proof}
補題 \ref{dpa-lemma-kernel} から直ちに従う。
\end{proof}

\begin{lemma}
\label{dpa-lemma-extend-to-completion}
$(A, I, \gamma)$ を分割冪環、$p$ を素数とする。$p$ が $A/I$ において
冪零ならば、次が成り立つ。
\begin{enumerate}
\item $p$-進完備化 $A^\wedge = \lim_e A/p^eA$ は $A/I$ へ全射する。
\item この写像の核は $I$ の $p$-進完備化 $I^\wedge$ である。
\item 各 $\gamma_n$ は $p$-進位相に関して連続であり、写像
$\gamma_n^\wedge : I^\wedge \to I^\wedge$ に延長され、$I^\wedge$ 上の
分割冪構造を定める。
\end{enumerate}
さらに $A$ が $\mathbf{Z}_{(p)}$-代数であるならば、
\begin{enumerate}
\item[(4)] 十分大きい $e$ に対して、イデアル $p^eA \subset I$ は
分割冪構造 $\gamma$ で保たれ、
$$
(A^\wedge, I^\wedge, \gamma^\wedge) = \lim_e (A/p^eA, I/p^eA, \bar\gamma)
$$
が分割冪環の圏で成り立つ。
\end{enumerate}
\end{lemma}

\begin{proof}
$p^tA/I=0$、すなわち $p^tA\subset I$ を満たす整数 $t\geq 1$ をとる。
写像 $A^\wedge\to A/I$ は合成
$A^\wedge\to A/p^tA\to A/I$ であり、これは全射である（たとえば
Algebra, Lemma \ref{algebra-lemma-completion-generalities} による）。
$e\geq t$ に対して
$p^eI \subset p^eA \cap I \subset p^{e - t}I$ であるから、合成
$A^\wedge\to A/I$ の核は $I$ の $p$-進完備化である。
写像 $\gamma_n$ は連続である。実際、
$$
\gamma_n(x + p^ey) =
\sum\nolimits_{i + j = n} p^{je}\gamma_i(x)\gamma_j(y) =
\gamma_n(x) \bmod p^eI
$$
が分割冪構造の公理から従う。写像 $\gamma_n^\wedge$ が $\gamma_n$ から
分割冪構造の公理を受け継ぐことは明らかである。最後の主張を示すため
$e>t$ とする。このとき $p^eA\subset I$ かつ
$\gamma_1(p^eA)\subset p^eA$ であり、$n>1$ に対して
$$
\gamma_n(p^ea) = p^n \gamma_n(p^{e - 1}a) = \frac{p^n}{n!} p^{n(e - 1)}a^n
\in p^e A
$$
である。これは $p^n/n! \in \mathbf{Z}_{(p)}$ であり、また
$n\geq 2$、$e\geq 2$ より $n(e-1)\geq e$ であるためである。
したがって $\gamma$ は $A/p^eA$ に延長される。補題
\ref{dpa-lemma-kernel} を参照せよ。極限に関する主張は、補題
\ref{dpa-lemma-limits} の証明における極限の構成から明らかである。
\end{proof}




\section{分割冪多項式代数}
\label{dpa-section-divided-power-polynomial-ring}

\noindent
非常に有用な例として、{\it 分割冪多項式代数}がある。
$A$ を環とし、$t \geq 1$ とする。次の $A$-代数を
$A\langle x_1, \ldots, x_t \rangle$ と表す。すなわち、
$A$-加群として
$$
A\langle x_1, \ldots, x_t \rangle =
\bigoplus\nolimits_{n_1, \ldots, n_t \geq 0} A x_1^{[n_1]} \ldots x_t^{[n_t]}
$$
とおき、乗法を
$$
x_i^{[n]}x_i^{[m]} = \frac{(n + m)!}{n!m!}x_i^{[n + m]}.
$$
で定める。また $x_i = x_i^{[1]}$ とおく。
$1 = x_1^{[0]} \ldots x_t^{[0]}$ であることに注意する。
同様の構成により、無限個の変数に関する分割冪多項式代数も得られる。
$n > 0$ に対して $x_i^{[n]}$ を零に写す標準的な $A$-代数準同型
$A\langle x_1, \ldots, x_t \rangle \to A$ が存在する。この写像の核を
$A\langle x_1, \ldots, x_t \rangle_{+}$ と表す。

\begin{lemma}
\label{dpa-lemma-divided-power-polynomial-algebra}
$(A, I, \gamma)$ を分割冪環とする。次のイデアル
$$
J = IA\langle x_1, \ldots, x_t \rangle + A\langle x_1, \ldots, x_t \rangle_{+}
$$
上には、次の条件を満たす分割冪構造 $\delta$ が一意に存在する。
\begin{enumerate}
\item $\delta_n(x_i) = x_i^{[n]}$ である。
\item $(A, I, \gamma) \to (A\langle x_1, \ldots, x_t \rangle, J, \delta)$
は分割冪環の準同型である。
\end{enumerate}
さらに、$(A\langle x_1, \ldots, x_t \rangle, J, \delta)$ は次の普遍性をもつ。
分割冪環の準同型
$\varphi : (A\langle x_1, \ldots, x_t \rangle, J, \delta) \to
(C, K, \epsilon)$
を与えることは、分割冪環の準同型 $A \to C$ と元
$k_1, \ldots, k_t \in K$ を与えることと同値である。
\end{lemma}

\begin{proof}
この補題を一変数の分割冪多項式代数の場合に証明する。
一般の場合も全く同様に論じられる。あるいは、
$A\langle x_1, \ldots, x_t\rangle$ が分割冪変数
$x_1, \ldots, x_t$ を一つずつ添加して得られる環と同型であることを用いてもよい。

\medskip\noindent
$A\langle x \rangle_{+}$ を
$x, x^{[2]}, x^{[3]}, \ldots$ で生成されるイデアルとする。
$J = IA\langle x \rangle + A\langle x \rangle_{+}$ であり、さらに
$$
IA\langle x \rangle \cap A\langle x \rangle_{+} =
IA\langle x \rangle \cdot A\langle x \rangle_{+}
$$
である。したがって、補題 \ref{dpa-lemma-two-ideals} により、イデアル
$IA\langle x \rangle$ と $A\langle x \rangle_{+}$ 上に分割冪構造が存在することを示せば十分である。
$A \to A\langle x \rangle$ は平坦なので、前者の存在は補題
\ref{dpa-lemma-gamma-extends} から従う。後者については、$A$ が無捩れならば、
補題 \ref{dpa-lemma-silly} (4) を適用して $\delta$ の存在が従う。
実際、元 $x^{[m]}$ を生成元として選ぶと、
$(x^{[m]})^n = \frac{(nm)!}{(m!)^n} x^{[nm]}$
であり、$n!$ は整数 $\frac{(nm)!}{(m!)^n}$ を割り切る。
一般の場合には、$A = R/\mathfrak a$ と書く。ただし $R$ はある無捩れ環
（例えば $\mathbf{Z}$ 上の多項式環）である。写像
$R\langle x \rangle \to A\langle x \rangle$ の核は
$\bigoplus \mathfrak a x^{[m]}$ である。補題 \ref{dpa-lemma-kernel} の
判定条件 (2)(c) を適用すれば、$R\langle x \rangle_{+}$ 上の分割冪構造が
所望の仕方で $A\langle x \rangle$ に延長されることが分かる。

\medskip\noindent
普遍性を証明する。分割冪環の準同型 $\varphi : A \to C$ と
$k_1, \ldots, k_t \in K$ が与えられたとき、
$$
A\langle x_1, \ldots, x_t \rangle \to C,\quad
x_1^{[n_1]} \ldots x_t^{[n_t]} \longmapsto
\epsilon_{n_1}(k_1) \ldots \epsilon_{n_t}(k_t)
$$
を考え、係数には $\varphi$ を用いる。これが $A$-代数準同型であることだけを
確かめればよい（詳細は省略する）。逆の構成は明らかである。
\end{proof}

\begin{remark}
\label{dpa-remark-divided-power-polynomial-algebra}
$(A, I, \gamma)$ を分割冪環とする。補題
\ref{dpa-lemma-divided-power-polynomial-algebra} には無限個の変数に関する変形版がある。
まず、$s < t$ ならば標準的な写像
$$
A\langle x_1, \ldots, x_s \rangle \to A\langle x_1, \ldots, x_t\rangle
$$
が存在することに注意する。したがって、任意の集合 $W$ に対して
$$
A\langle x_w: w \in W\rangle =
\colim_{E \subset W} A\langle x_e:e \in E\rangle
$$
とおく（余極限は $W$ の有限部分集合 $E$ 全体にわたり、遷移写像は上記のものである）。
余極限の定義から、$A\langle x_w : w \in W\rangle$ の普遍写像性は、
補題 \ref{dpa-lemma-divided-power-polynomial-algebra} に述べた写像性と
完全に同様であることが分かる。
\end{remark}

\noindent
次の補題は \cite{BO} に見いだされる。

\begin{lemma}
\label{dpa-lemma-need-only-gamma-p}
$p$ を素数とする。$A$ を、$p$ で割り切れない任意の整数 $n$ が可逆である環、
すなわち $\mathbf{Z}_{(p)}$-代数とする。$I \subset A$ をイデアルとする。
$I$ 上の二つの分割冪構造 $\gamma, \gamma'$ が等しいための必要十分条件は
$\gamma_p = \gamma'_p$ である。さらに、写像 $\delta : I \to I$ が
\begin{enumerate}
\item 任意の $x \in I$ に対して $p!\delta(x) = x^p$ である\footnote{この条件は
他の二条件から従う。注意 \ref{dpa-remark-ryo-suzuki} を参照せよ。}。
\item 任意の $a \in A$、$x \in I$ に対して $\delta(ax) = a^p\delta(x)$ である。
\item
$\delta(x + y) =
\delta(x) +
\sum\nolimits_{i + j = p, i,j \geq 1} \frac{1}{i!j!} x^i y^j +
\delta(y)$ が任意の $x, y \in I$ に対して成り立つ。
\end{enumerate}
を満たすならば、$\gamma_p = \delta$ となる $I$ 上の分割冪構造 $\gamma$ が一意に存在する。
\end{lemma}

\begin{proof}
$n$ が $p$ で割り切れないならば、
$\gamma_n(x) = c x\gamma_{n - 1}(x)$ である。ただし $c$ は
$\mathbf{Z}_{(p)}$ の単元である。さらに、
$$
\gamma_{pm}(x) = c \gamma_m(\gamma_p(x))
$$
であり、ここでも $c$ は $\mathbf{Z}_{(p)}$ の単元である。したがって第一の主張は明らかである。
第二の主張については、これらの等式を逆向きに用いてすべての $\gamma_n$ を定義できる。
より正確には、$n = a_0 + a_1p + \ldots + a_e p^e$、
$a_i \in \{0, \ldots, p - 1\}$ ならば、
$$
\gamma_n(x) = c_n x^{a_0} \delta(x)^{a_1} \ldots \delta^e(x)^{a_e}
$$
とおく。ただし $c_n \in \mathbf{Z}_{(p)}$ は
$$
c_n =
{(p!)^{a_1 + a_2(1 + p) + \ldots + a_e(1 + \ldots + p^{e - 1})}}/{n!}.
$$
で定義される。ここで、定義 \ref{dpa-definition-divided-powers} にある
分割冪構造の公理 (1)--(5) を示さなければならない。
(1) と (3) は直ちに従う。(2) と (5) は直接計算で確かめられるので省略する。
$x, y \in I$ とする。環準同型
$$
\varphi : \mathbf{Z}_{(p)}\langle u, v \rangle \longrightarrow A
$$
であって、$u^{[n]}$ を $\gamma_n(x)$ に、$v^{[n]}$ を $\gamma_n(y)$ に写すものが
存在すると主張する。$\mathbf{Z}_{(p)}\langle u, v\rangle$ の構成から、これは
$$
\gamma_n(x)\gamma_m(x) = \frac{(n + m)!}{n!m!} \gamma_{n + m}(x)
$$
が $A$ において成り立ち、$y$ についても同様であることを確かめるということである。
これは $\gamma$ が (2) を満たすことから従う。
$\epsilon$ を、$\mathbf{Z}_{(p)}\langle u, v\rangle$ のイデアル
$\mathbf{Z}_{(p)}\langle u, v\rangle_{+}$ 上の分割冪構造とする。
次に、$f \in \mathbf{Z}_{(p)}\langle u, v\rangle_{+}$ と任意の $n$ に対して
$\varphi(\epsilon_n(f)) = \gamma_n(\varphi(f))$ であると主張する。
$n = 0, 1, \ldots, p - 1$ では明らかである。$n = p$ の場合には、
イデアル $\mathbf{Z}_{(p)}\langle u, v\rangle_{+}$ の一組の生成元について示せば十分である。
実際、$\epsilon_p$ と $\gamma_p = \delta$ はともに補題の性質 (1) と (3) を満たす。
したがって、
$\gamma_p(\gamma_n(x)) = \frac{(pn)!}{p!(n!)^p}\gamma_{pn}(x)$
および $y$ に対する同様の等式を示せばよいが、これは $\gamma$ が (5) を満たすことから従う。
さて、$n = a_0 + a_1p + \ldots + a_e p^e$ が上のように $p$ 進展開された任意の整数ならば、
$$
\epsilon_n(f) =
c_n f^{a_0} \gamma_p(f)^{a_1} \ldots \gamma_p^e(f)^{a_e}
$$
である。これは $\epsilon$ が分割冪構造だからである。したがって、任意の $n$ に対して
$\varphi(\epsilon_n(f)) = \gamma_n(\varphi(f))$ が成り立つ。
これを $f = u + v$ に適用すれば、$\epsilon$ が分割冪構造であることから
$\gamma$ が公理 (4) を満たすことが分かる。
\end{proof}

\begin{remark}
\label{dpa-remark-ryo-suzuki}
$A \supset I$ を補題 \ref{dpa-lemma-need-only-gamma-p} のとおりとし、
$\delta : I \to I$ を補題の (2) と (3) を満たす写像とする。このとき (1) も成り立つ。
まず $n$ に関する帰納法により、$n \geq 1$ に対して
$\delta(nx) = n\delta(x) + \frac{n^p - n}{p!}x^p$ であることを示す。
$n = 1$ の場合は成り立つ。ある $n$ に対して成り立つと仮定すると、
$\delta((n + 1)x) = \delta(nx) + \delta(x) +
(\sum_{i = 1}^{p - 1} \frac{n^i}{i!(p - i)!})x^p
= (n + 1)\delta(x) + \left( \frac{n^p - n}{p!} + \frac{(n + 1)^p}{p!}
- \frac{n^p + 1}{p!} \right)x^p =
(n + 1)\delta(x) + \frac{(n + 1)^p - (n + 1)}{p!}x^p$,
ゆえに $n + 1$ に対しても成り立つ。一方、$\delta(nx) = n^p\delta(x)$ であるから、
$(n^p - n)\delta(x) = \frac{n^p - n}{p!}x^p$. 
任意の $p$ に対して、$p^2 \not \mid n^p - n$ となる $n$ が存在する。
実際、
$(\mathbf{Z}/p^2\mathbf{Z})^\times \cong \mathbf{Z}/p(p-1)\mathbf{Z}$.
したがって、$n^{p-1} \not\equiv 1 \mod p^2$ となる $n$ が存在する。
ゆえに $p!\delta(x) = x^p$ を得る。
\end{remark}









\section{Tate 分解}
\label{dpa-section-tate}

\noindent
本節では、分割冪構造と微分次数付き代数を組み合わせた
\cite{Tate-homology} および \cite{AH} の分解について簡単に論じる。
本節では微分次数付き代数に{\it ホモロジー記法}を用いる。
考える微分次数付き代数は非負のホモロジー次数に置かれる。
したがって、微分次数付き代数 $(A, \text{d})$ は乗法を備えた鎖複体
$$
\ldots \to A_2 \to A_1 \to A_0 \to 0 \to \ldots
$$
として与えられる。

\medskip\noindent
$R$ を環（いつものように可換環）とする。本節ではしばしば、負次数の成分が零である
次数付き $R$-代数 $A = \bigoplus_{d \geq 0} A_d$ を考える。
$A_+ = \bigoplus_{d > 0} A_d$ とおく。また、
$A_{even} = \bigoplus_{d \geq 0} A_{2d}$ および
$A_{odd} = \bigoplus_{d \geq 0} A_{2d + 1}$ と書く。
斉次元 $x, y$ に対して
$xy = (-1)^{\deg(x)\deg(y)}yx$ となるとき、$A$ は次数可換であるという。
さらに奇数次数の斉次元 $x$ に対して $x^2 = 0$ となるとき、
$A$ は厳密次数可換であるという。最後に、次の定義を理解するため、
$A$ が $\mathbf{Q}$-代数ならば $\gamma_n(x) = x^n/n!$ であることを念頭に置く。

\begin{definition}
\label{dpa-definition-divided-powers-graded}
$R$ を環とする。$A = \bigoplus_{d \geq 0} A_d$ を厳密次数可換な
次数付き $R$-代数とする。すべての $n > 0$ に対して定義された写像の族
$\gamma_n : A_{even, +} \to A_{even, +}$ が次を満たすとき、これを
$A$ 上の{\it 分割冪構造}という。
\begin{enumerate}
\item $x \in A_{2d}$ ならば $\gamma_n(x) \in A_{2nd}$ である。
\item 任意の $x$ に対して $\gamma_1(x) = x$ であり、また $\gamma_0(x) = 1$ とおく。
\item $\gamma_n(x)\gamma_m(x) = \frac{(n + m)!}{n! m!} \gamma_{n + m}(x)$,
\item すべての $x \in A_{even}$ と $y \in A_{even, +}$ に対して
$\gamma_n(xy) = x^n\gamma_n(y)$ である。
\item $x, y \in A_{odd}$ が斉次で $n > 1$ ならば $\gamma_n(xy) = 0$ である。
\item $x, y \in A_{even, +}$ ならば
$\gamma_n(x + y) = \sum_{i = 0, \ldots, n} \gamma_i(x)\gamma_{n - i}(y)$ である。
\item $\gamma_n(\gamma_m(x)) =
\frac{(nm)!}{n! (m!)^n} \gamma_{nm}(x)$ が $x \in A_{even, +}$ に対して成り立つ。
\end{enumerate}
\end{definition}

\noindent
条件 (2)、(3)、(4)、(6)、(7) から、$\gamma$ は可換環
$A_{even}$ のイデアル $A_{even, +}$ 上の「通常の」分割冪構造であることが分かる。
第 \ref{dpa-section-divided-powers} 節、
第 \ref{dpa-section-divided-power-rings} 節、
第 \ref{dpa-section-extend} 節、および
第 \ref{dpa-section-divided-power-polynomial-ring} 節を参照せよ。
特に、すべての $x \in A_{even, +}$ に対して $n!\gamma_n(x) = x^n$ である。
条件 (1) は $\gamma$ が次数付けと両立することを述べ、条件 (5) は $n > 1$ のとき
$\gamma_n$ が奇数次数の斉次元の積上で零になることを述べている。しかし、
$$
\gamma_2(z_1 z_2 + z_3 z_4) = z_1z_2z_3z_4
$$
が非零となることもあり得る。ここで $z_1, z_2, z_3, z_4$ は奇数次数の斉次元である。

\begin{example}[奇数次数の変数の添加]
\label{dpa-example-adjoining-odd}
$R$ を環とする。$(A, \gamma)$ を、上の定義における分割冪構造を備えた
厳密次数可換な次数付き $R$-代数とする。$d > 0$ を奇数とする。
この状況では、次数 $d$ の変数 $T$ を $A$ に添加できる。すなわち、
$$
A\langle T \rangle = A \oplus AT
$$
とおき、次数付けを $A\langle T \rangle_m = A_m \oplus A_{m - d}T$ で定める。
$A\langle T \rangle$ 上には、$A$ 上の与えられた分割冪構造と両立する
分割冪構造が一意に存在すると主張する。すなわち、
$x \in A_{even, +}$ および $y \in A_{odd}$ に対して
$$
\gamma_n(x + yT) = \gamma_n(x) + \gamma_{n - 1}(x)yT
$$
とおく。
\end{example}

\begin{example}[偶数次数の変数の添加]
\label{dpa-example-adjoining-even}
$R$ を環とする。$(A, \gamma)$ を、上の定義における分割冪構造を備えた
厳密次数可換な次数付き $R$-代数とする。$d > 0$ を偶数とする。
この状況では、次数 $d$ の変数 $T$ を $A$ に添加できる。すなわち、
$$
A\langle T \rangle = A \oplus AT \oplus AT^{(2)} \oplus AT^{(3)} \oplus \ldots
$$
とおき、乗法を
$$
T^{(n)} T^{(m)} = \frac{(n + m)!}{n!m!} T^{(n + m)}
$$
で定め、次数付けを
$$
A\langle T \rangle_m =
A_m \oplus A_{m - d}T \oplus A_{m - 2d}T^{(2)} \oplus \ldots
$$
で定める。$A\langle T \rangle$ 上には、$A$ 上の与えられた分割冪構造と両立し、
$\gamma_n(T^{(i)}) = T^{(ni)}$ を満たす分割冪構造が一意に存在すると主張する。
この分割冪構造を定義するため、まず $x_i \in A_{even}$ のとき
$$
\gamma_n\left(\sum\nolimits_{i > 0} x_i T^{(i)}\right) =
\sum \prod\nolimits_{n = \sum e_i} x_i^{e_i} T^{(ie_i)}
$$
とおく。$x_0 \in A_{even, +}$ ならば、
$$
\gamma_n\left(\sum\nolimits_{i \geq 0} x_i T^{(i)}\right) =
\sum\nolimits_{a + b = n}
\gamma_a(x_0)\gamma_b\left(\sum\nolimits_{i > 0} x_iT^{(i)}\right)
$$
とおく。ただし $\gamma_b$ は上で定義したものである。
\end{example}

\begin{remark}
\label{dpa-remark-adjoining-set-of-variables}
例 \ref{dpa-example-adjoining-odd} および \ref{dpa-example-adjoining-even} のように
一つだけでなく、与えられた次数 $d > 0$ の外積生成元または分割冪生成元の集合
（無限集合でもよい）を添加することもできる。すなわち、注意
\ref{dpa-remark-divided-power-polynomial-algebra} に従い、上の $(A, \gamma)$ と集合 $J$ に対して、
$A\langle T_j : j \in J\rangle$ を、$J$ の有限部分集合 $S$ 全体にわたる代数
$A\langle T_j : j \in S\rangle$ の有向余極限とする。
この代数には、$A$ 上の与えられた構造および各生成元 $T_j$ 上の構造と両立する
分割冪構造が一意に存在することは直ちに分かる。
\end{remark}

\noindent
ここで分割冪構造の定義を微分と結び付ける。次の定義を理解するため、
$A$ が $\mathbf{Q}$-代数で $x \in A_{even, +}$ ならば
$\text{d}(x^n/n!) = \text{d}(x)x^{n - 1}/(n - 1)!$ であることに注意する。

\begin{definition}
\label{dpa-definition-divided-powers-dga}
$R$ を環とする。$A = \bigoplus_{d \geq 0} A_d$ を厳密次数可換な
微分次数付き $R$-代数とする。$A$ 上の分割冪構造 $\gamma$ が、
すべての $x \in A_{even, +}$ に対して
$\text{d}(\gamma_n(x)) = \text{d}(x)\gamma_{n - 1}(x)$ を満たすとき、
$\gamma$ は{\it 微分次数付き構造と両立する}という。
\end{definition}

\noindent
注意：$(A, \text{d}, \gamma)$ を定義
\ref{dpa-definition-divided-powers-dga} のとおりとする。
$x$ が境界であっても、$\gamma_n(x)$ が境界であるとは限らない。
したがって、一般には $\gamma$ はホモロジー代数 $H(A)$ 上に分割冪構造を誘導しない。
いくつかの論文では、これを保証するために著者が追加の両立条件を課しているが、
ここではそうしないことを選ぶ。

\begin{lemma}
\label{dpa-lemma-dpdga-good}
$(A, \text{d}, \gamma)$ および $(B, \text{d}, \gamma)$ を定義
\ref{dpa-definition-divided-powers-dga} のとおりとする。
$f : A \to B$ を分割冪構造と両立する微分次数付き代数の写像とする。次を仮定する。
\begin{enumerate}
\item $k > 0$ に対して $H_k(A) = 0$ である。
\item $f$ は全射である。
\end{enumerate}
このとき $\gamma$ は次数付き $R$-代数 $H(B)$ 上に分割冪構造を誘導する。
\end{lemma}

\begin{proof}
$x$ と $x'$ を同じ次数 $2d$ の斉次元で、$H(B)$ において同じコホモロジー類を
定めるものとする。$x' - x = \text{d}(w)$ と書く。
$x$ の持ち上げ $y \in A_{2d}$ と $w$ の持ち上げ $z \in A_{2d + 1}$ を選ぶ。
このとき $y' = y + \text{d}(z)$ は $x'$ の持ち上げである。したがって、
$$
\gamma_n(y') = \sum \gamma_i(y) \gamma_{n - i}(\text{d}(z))
= \gamma_n(y) +
\sum\nolimits_{i < n} \gamma_i(y) \gamma_{n - i}(\text{d}(z))
$$
$A$ は正次数で非輪状であり、すべての $j$ に対して
$\text{d}(\gamma_j(\text{d}(z))) = 0$ であるから、これは
$$
\gamma_n(y') = \gamma_n(y) +
\sum\nolimits_{i < n} \gamma_i(y) \text{d}(z_i)
$$
と書ける。ただし $z_i$ は $A$ のある元である。さらに、$0 < i < n$ に対して
$$
\text{d}(\gamma_i(y) z_i) =
\text{d}(\gamma_i(y))z_i + \gamma_i(y)\text{d}(z_i) =
\text{d}(y) \gamma_{i - 1}(y) z_i + \gamma_i(y)\text{d}(z_i)
$$
であり、第一項は $B$ において零に写る。実際、$\text{d}(y)$ は $B$ において零に写る。
したがって、$\gamma_n(x')$ と $\gamma_n(x)$ は $H(B)$ の同じ元に写る。
こうして、すべての $d > 0$ と $n > 0$ に対して良定義な写像
$\gamma_n : H_{2d}(B) \to H_{2nd}(B)$ が得られる。
これが $H(B)$ 上の分割冪構造を定めることの確認は省略する。
\end{proof}

\begin{lemma}
\label{dpa-lemma-base-change-div}
$(A, \text{d}, \gamma)$ を定義
\ref{dpa-definition-divided-powers-dga} のとおりとする。$R \to R'$ を環準同型とする。
このとき $\text{d}$ と $\gamma$ は $A' = A \otimes_R R'$ 上に同様の構造を誘導し、
$(A', \text{d}, \gamma)$ は定義 \ref{dpa-definition-divided-powers-dga} のとおりとなる。
\end{lemma}

\begin{proof}
$A'_{even} = A_{even} \otimes_R R'$ および
$A'_{even, +} = A_{even, +} \otimes_R R'$ であることに注意する。
したがって示すべきことは、分割冪 $\gamma$ が $A'_{even}$ に延長されることである
（用語は定義 \ref{dpa-definition-extends} のとおり）。$\gamma$ が延長されることを示せば、
この延長が所望の性質をすべてもつことは容易に従う。

\medskip\noindent
多項式 $R$-代数 $P$（生成元は任意の集合でよい）と全射な $R$-代数準同型
$P \to R'$ を選ぶ。環準同型 $A_{even} \to A_{even} \otimes_R P$ は平坦であるから、
補題 \ref{dpa-lemma-gamma-extends} により、分割冪 $\gamma$ は
$A_{even} \otimes_R P$ に一意に延長される。$J = \Ker(P \to R')$ とおく。
$\gamma$ が $A \otimes_R R'$ に延長されることを示すには、
$I' = \Ker(A_{even, +} \otimes_R P \to A_{even, +} \otimes_R R')$ が、
すべての $n > 0$ に対して $\gamma_n(z) \in I'$ となる元 $z$ で生成されることを
示せば十分である。これは明らかである。実際、$I'$ は
$x \in A_{even, +}$ および $f \in \Ker(P \to R')$ による
$x \otimes f$ の形の元で生成される。
\end{proof}

\begin{lemma}
\label{dpa-lemma-extend-differential}
$(A, \text{d}, \gamma)$ を定義 \ref{dpa-definition-divided-powers-dga} のとおりとする。
$d \geq 1$ を整数とする。$A\langle T \rangle$ を、例
\ref{dpa-example-adjoining-odd} または \ref{dpa-example-adjoining-even} で構成した、
$\deg(T) = d$ である $T$ に関する次数付き分割冪多項式代数とする。
$f \in A_{d - 1}$ を $\text{d}(f) = 0$ となる元とする。
$A\langle T \rangle$ 上には、$\text{d}(T) = f$ を満たし、かつ
$A\langle T \rangle$ 上の分割冪構造と両立する微分 $\text{d}$ が一意に存在する。
\end{lemma}

\begin{proof}
省略する直接計算によって証明される。
\end{proof}

\noindent
補題 \ref{dpa-lemma-compute-cohomology-adjoin-variable} では、いくつかの特別な場合に
$A\langle T \rangle$ のコホモロジーを計算する。
以下が Tate の構成であり、Avramov と Halperin によって拡張されたものである。

\begin{lemma}
\label{dpa-lemma-tate-resolution}
$R \to S$ を可換環の準同型とする。次の性質をもつ分解
$$
R \to A \to S
$$
が存在する。
\begin{enumerate}
\item $(A, \text{d}, \gamma)$ は定義
\ref{dpa-definition-divided-powers-dga} のとおりである。
\item $S$ に零微分を入れると、$A \to S$ は擬同型である。
\item $A_0 = R[x_j : j \in J] \to S$ は、多項式環から $S$ への任意の全射である。
\item $A$ は $R$ 上の次数付き分割冪多項式代数である。
\end{enumerate}
最後の条件は、例 \ref{dpa-example-adjoining-odd} または
\ref{dpa-example-adjoining-even} のように、各正次数で変数 $T$ の集合を
順次添加することにより $A_0$ から $A$ が構成されることを意味する。
さらに、$R$ が Noether 環で $R \to S$ が有限型ならば、各次数で生成元が
有限個だけとなるように $A$ を選べる。
\end{lemma}

\begin{proof}
$R$ が Noether 環で $R \to S$ が有限型の場合について構成を詳述する。
これらの仮定がない場合も証明は同じであるが、各次数で生成元の集合
（無限集合でもよい）を用いなければならない。

\medskip\noindent
構成の開始：$A(0) = R[x_1, \ldots, x_n]$ を通常の多項式環とし、
$A(0) \to S$ を全射とする。次数付けとして
$A(0)_0 = A(0)$、および $d \not = 0$ に対して $A(0)_d = 0$ とおく。
したがって $\text{d} = 0$ であり、$n > 0$ に対する $\gamma_n$ も零である。

\medskip\noindent
与えられた写像 $A(0) = R[x_1, \ldots, x_n] \to S$ の核の生成元
$f_1, \ldots, f_m \in R[x_1, \ldots, x_n]$ を選ぶ。
例 \ref{dpa-example-adjoining-odd} を $m$ 回適用して、次数付き分割冪多項式代数
$$
A(1) = A(0)\langle T_1, \ldots, T_m\rangle
$$
を得る。ただし $\deg(T_i) = 1$ である。$\text{d}(T_i) = f_i$ とおく。
$A(1)$ は $A(0)$ 上の分割冪多項式代数であり、$\text{d}(f_i) = 0$ なので、
補題 \ref{dpa-lemma-extend-differential} により、これは $A(1)$ 上の微分に一意に延長される。

\medskip\noindent
帰納法の仮定：分解
$$
R \to A(0) \to A(1) \to \ldots \to A(m) \to S
$$
が与えられていると仮定する。ここで $A(0)$ と $A(1)$ は上のとおりであり、
$2 \leq m' \leq m$ に対する各 $R \to A(m') \to S$ は補題の主張の性質 (1) と (4) を満たし、
(2) の代わりに、$m' > i \geq 0$ のとき
$H_i(A(m')) \to H_i(S)$ が同型であるという条件を満たす。基底の場合は $m = 1$ である。

\medskip\noindent
帰納段階：帰納法の仮定のとおりの $R \to A(m) \to S$ があるとする。
群 $H_m(A(m))$ を考える。これは $H_0(A(m)) = S$ 上の加群である。
実際、これは $A(m)_m$ の部分商であり、$A(m)_m$ は
$A(m)_0 = R[x_1, \ldots, x_n]$ 上の有限型加群である。
したがって、有限個の元
$$
e_1, \ldots, e_t \in \Ker(\text{d} : A(m)_m \to A(m)_{m - 1})
$$
で、その像がこの加群を生成するものを選べる。例 \ref{dpa-example-adjoining-odd} または
\ref{dpa-example-adjoining-even} を $t$ 回適用して、次数付き分割冪代数
$$
A(m + 1) = A(m)\langle T_1, \ldots, T_t\rangle
$$
を得る。ただし $\deg(T_i) = m + 1$ である。$\text{d}(T_i) = e_i$ とおく。
$A(m + 1)$ は $A(m)$ 上の分割冪多項式代数であり、$\text{d}(e_i) = 0$ なので、
これは分割冪構造と両立する $A(m + 1)$ 上の微分に一意に延長される。
次数 $m + 1$ 以上にだけ素材を添加したので、$i < m$ に対して
$H_i(A(m + 1)) = H_i(A(m))$ である。さらに、構成から
$H_m(A(m + 1)) = 0$ であることは明らかである。

\medskip\noindent
証明を終えるため、写像の列
$$
R \to A(0) \to A(1) \to \ldots \to A(m) \to A(m + 1) \to \ldots \to S
$$
が存在することを示した点に注意する。$A = \colim A(m)$ とおけば証明が完了する。
\end{proof}

\begin{lemma}
\label{dpa-lemma-tate-resoluton-pseudo-coherent-ring-map}
$R \to S$ を擬コヒーレントな環準同型とする（More on Algebra, Definition
\ref{more-algebra-definition-pseudo-coherent-perfect}）。このとき補題
\ref{dpa-lemma-tate-resolution} が成り立ち、$S$ の分解 $A$ は各次数で
有限個の生成元だけをもつように選べる。
\end{lemma}

\begin{proof}
補題 \ref{dpa-lemma-tate-resolution} と全く同じ方法で証明される。
唯一の追加点は、$A(m) \to S$ が与えられたとき、
$H_m = H_m(A(m))$ が有限 $R[x_1, \ldots, x_m]$-加群であることを
示さなければならないことである（そうすれば次の段階で有限個の変数だけを添加すればよい）。
複体
$$
\ldots \to A(m)_{m - 1} \to A(m)_m \to A(m)_{m - 1} \to
\ldots \to A(m)_0 \to S \to 0
$$
を考える。$S$ は擬コヒーレントな $R[x_1, \ldots, x_n]$-加群であり、
$A(m)_i$ は有限自由 $R[x_1, \ldots, x_n]$-加群なので、これは擬コヒーレントな複体である。
More on Algebra, Lemma \ref{more-algebra-lemma-complex-pseudo-coherent-modules} を参照せよ。
この複体はホモロジー次数 $> m$ で完全なので、More on Algebra, Lemma
\ref{more-algebra-lemma-finite-cohomology} により $H_m$ は有限 $R$-加群である。
\end{proof}

\begin{lemma}
\label{dpa-lemma-uniqueness-tate-resolution}
$R$ を可換環とする。$(A, \text{d}, \gamma)$ および
$(B, \text{d}, \gamma)$ を定義 \ref{dpa-definition-divided-powers-dga} のとおりとする。
$\overline{\varphi} : H_0(A) \to H_0(B)$ を $R$-代数準同型とする。次を仮定する。
\begin{enumerate}
\item $A$ は $R$ 上の次数付き分割冪多項式代数である。
\item $k > 0$ に対して $H_k(B) = 0$ である。
\end{enumerate}
このとき、$\overline{\varphi}$ を持ち上げる、分割冪と両立する微分次数付き
$R$-代数の写像 $\varphi : A \to B$ が存在する。
\end{lemma}

\begin{proof}
この仮定は、次数零の多項式生成元の集合、正の奇数次数の外積生成元、
および正の偶数次数の分割冪生成元を $R$ に順次添加することで $A$ が得られることを意味する。
したがって、$A$ にはフィルトレーション
$R \subset A(0) \subset A(1) \subset \ldots$ があり、$A(m + 1)$ は
$A(m)$ に次数 $m + 1$ の適切な型の生成元（これらを単に「分割冪生成元」と呼ぶ）を
添加して得られる。特に、$A(0) \to H_0(A)$ は、$R$ 上の通常の多項式代数から
$H_0(A)$ への全射である。したがって、$\overline{\varphi}$ を
$R$-代数準同型 $\varphi(0) : A(0) \to B_0$ に持ち上げられる。

\medskip\noindent
$A(1) = A(0)\langle T_j : j \in J\rangle$ と書く。ここで $J$ は次数 $1$ の
分割冪変数 $T_j$ のある集合である。
$f_j = \varphi(0)(\text{d}(T_j)) \in B_0$ とおく。
$\text{d}T_j$ は $H_0(A)$ において零に写るので、$f_j$ は $H_0(B)$ において零に写る。
したがって、$\text{d}(b_j) = f_j$ となる $b_j \in B_1$ が見つかる。
補題 \ref{dpa-lemma-divided-power-polynomial-algebra} の分割冪多項式代数の普遍性により、
$T_j$ を $f_j$ に写す $\varphi(0)$ の持ち上げ
$\varphi(1) : A(1) \to B$ が得られる。

\medskip\noindent
ある $m \geq 1$ に対して $\varphi(m)$ を構成したならば、全く同じ方法で
$\varphi(m + 1) : A(m + 1) \to B$ を構成できる。詳細は省略する。
\end{proof}

\begin{lemma}
\label{dpa-lemma-divided-powers-on-tor}
$R$ を可換環とし、$S$ と $T$ を可換 $R$-代数とする。このとき
$$
\operatorname{Tor}_*^R(S, T).
$$
上には、分割冪を備えた厳密次数可換 $R$-代数の標準的な構造が存在する。
\end{lemma}

\begin{proof}
上のような分解 $R \to A \to S$ を選ぶ。$A \to S$ は擬同型であり、
$A_d$ は自由 $R$-加群なので、微分次数付き代数 $B = A \otimes_R T$ は
補題に表示した Tor 群を計算する。全射 $R[y_j : j \in J] \to T$ を選ぶ。
すると $B$ は微分次数付き代数 $A[y_j : j \in J]$ の商であり、後者の
ホモロジーは次数 $0$ に集中する（それは $S[y_j : j \in J]$ に等しい）。
補題 \ref{dpa-lemma-base-change-div} により、微分次数付き代数 $B$ および
$A[y_j : j \in J]$ は、微分と両立する分割冪構造をもつ。
したがって補題 \ref{dpa-lemma-dpdga-good} により、$H(B)$ 上に所望の分割冪構造を得る。

\medskip\noindent
このように構成した分割冪代数構造は $A$ の選択によらない。
実際、$A'$ が第二の選択ならば、補題 \ref{dpa-lemma-uniqueness-tate-resolution} により、
すべての構造および $S$ への増大写像を保つ写像 $A \to A'$ が存在する。
すると誘導写像
$B = A \otimes_R T \to A' \otimes_R T' = B'$ もすべての構造を保ち、擬同型である。
したがって、誘導される Tor 代数の同型は積および分割冪と両立する。
\end{proof}





\section{完全交叉への応用}
\label{dpa-section-application-ci}

\noindent
$R$ を環とし、$(A, \text{d}, \gamma)$ を定義
\ref{dpa-definition-divided-powers-dga} のとおりとする。次数 $2$ の{\it 導分}とは、
次の性質をもつ $R$-線形写像 $\theta : A \to A$ のことである。
\begin{enumerate}
\item $\theta(A_d) \subset A_{d - 2}$,
\item $\theta(xy) = \theta(x)y + x\theta(y)$,
\item $\theta$ は $\text{d}$ と可換である。
\item $\theta(\gamma_n(x)) = \theta(x) \gamma_{n - 1}(x)$
がすべての $d$ および $x \in A_{2d}$ に対して成り立つ。
\end{enumerate}
次の補題で導分を構成する。

\begin{lemma}
\label{dpa-lemma-get-derivation}
$R$ を環とし、$(A, \text{d}, \gamma)$ を定義
\ref{dpa-definition-divided-powers-dga} のとおりとする。$R' \to R$ を、その核が
平方零で一つの元 $f$ によって生成される環の全射とする。
$A$ が $R$ 上の次数付き分割冪多項式代数で、各次数の変数が有限個ならば、
導分 $\theta : A/IA \to A/IA$ が得られる。ここで $I$ は $R$ における
$f$ の零化イデアルである。
\end{lemma}

\begin{proof}
$A$ は分割冪多項式代数なので、$A = A' \otimes_{R'} R$ となる
$R'$ 上の分割冪多項式代数 $A'$ が見つかる。さらに、$\text{d}$ を
$A'$ 上の $R$-線形作用素 $\text{d}$ に持ち上げて、次を満たすようにできる。
\begin{enumerate}
\item  $\text{d}(xy) = \text{d}(x)y + (-1)^{\deg(x)}x \text{d}(y)$
が $A'$ の斉次元 $x, y$ に対して成り立つ。
\item  $\text{d}(\gamma_n(x)) = \text{d}(x) \gamma_{n - 1}(x)$ が
$x \in A'_{even, +}$ に対して成り立つ。
\end{enumerate}
詳細は省略する（ヒント：変数を一つずつ扱う）。しかし、$\text{d}^2$ は
$A'$ 上で零とは限らない。$\text{d}^2$ が $A'$ を
$fA' \cong A/IA$ に写すことは明らかである。したがって $\text{d}^2$ は
$fA'$ を零化し、写像 $A \to A/IA$ を経由する。$\text{d}^2$ は $R$-線形なので、
所望の写像 $\theta : A/IA \to A/IA$ を得る。導分の性質の確認は直ちにできる。
\end{proof}

\begin{lemma}
\label{dpa-lemma-compute-theta}
仮定と記法を補題 \ref{dpa-lemma-get-derivation} のとおりとする。
$S = H_0(A)$ が、ある $n,m$ および $f_j \in R[x_1, \ldots, x_n]$ に対して
$R[x_1, \ldots, x_n]/(f_1, \ldots, f_m)$
と同型であると仮定する。さらに、関係式
$$
\sum r_j f_j = 0
$$
が与えられているとする。ただし $r_j \in R[x_1, \ldots, x_n]$ である。
$r_j, f_j$ の持ち上げ $r'_j, f'_j \in R'[x_1, \ldots, x_n]$ を選ぶ。
ある $g \in R/I[x_1, \ldots, x_n]$ によって $\sum r'_j f'_j = gf$ と書く。
$H_1(A) = 0$ で、各 $r_j$ のすべての係数が $I$ に属するならば、
$S/IS$ において $\theta(\xi) = g$ となる元 $\xi \in H_2(A/IA)$ が存在する。
\end{lemma}

\begin{proof}
$A(0) \subset A(1) \subset A(2) \subset \ldots$ を、$A(m)$ が
$A(m - 1)$ に次数 $m$ の分割冪変数を添加して得られるような $A$ のフィルトレーションとする。
すると $A(0)$ は $R$ 上の多項式代数で、$R$-代数の全射 $A(0) \to S$ を備える。
したがって、$S$ への増大写像を持ち上げる写像
$$
\varphi : R[x_1, \ldots, x_n] \to A(0)
$$
を選べる。次に、次数 $1$ のある分割冪変数 $T_i$ に対して
$A(1) = A(0)\langle T_1, \ldots, T_t\rangle$ である。$H_0(A) = S$ なので、
$\text{d}(\xi_j) = \varphi(f_j)$ となる
$\xi_j \in \sum A(0)T_i$ を選べる。すると
$$
\text{d}\left(\sum \varphi(r_j) \xi_j\right) =
\sum  \varphi(r_j) \varphi(f_j) =  \sum \varphi(r_jf_j) = 0
$$
$H_1(A) = 0$ なので、$\text{d}(\xi) = \sum \varphi(r_j)\xi_j$ となる
$\xi \in A_2$ を選べる。$r_j$ の係数が $I$ に属するならば、
$\varphi(r_j)$ についても同じである。この場合 $\text{d}(\xi)$ は
$A_1/IA_1$ において消えるので、$\xi$ は $H_2(A/IA)$ の類を定める。

\medskip\noindent
補題 \ref{dpa-lemma-get-derivation} の証明における $\theta$ の構成では、
$A(i)$ を $A'(i)$ に順次持ち上げ、微分 $\text{d}$ も持ち上げる。
$\varphi$ を $\varphi' : R'[x_1, \ldots, x_n] \to A'(0)$ に持ち上げる。
次に $A'(1) = A'(0)\langle T_1, \ldots, T_t\rangle$ である。
さらに、$\xi_j$ を $\xi'_j \in \sum A'(0)T_i$ に持ち上げられる。
すると、ある $a_j \in A'(0)$ に対して
$\text{d}(\xi'_j) = \varphi'(f'_j) + f a_j$ である。
$\xi$ の持ち上げ $\xi' \in A'_2$ を考える。このとき
$$
\text{d}(\xi') = \sum \varphi'(r'_j)\xi'_j + \sum fb_iT_i
$$
である。ただし $b_i \in A(0)$ である。もう一度 $\text{d}$ を適用すると、
$$
\theta(\xi) = \sum \varphi'(r'_j)\varphi'(f'_j) +
\sum f \varphi'(r'_j) a_j + \sum fb_i \text{d}(T_i)
$$
第一項が所望の結果を与える。第二項は零である。実際、$r_j$ の係数は $I$ に属し、
したがって $f$ により零化される。第三項は $H_0$ において零に写る。
実際、$\text{d}(T_i)$ は零に写る。
\end{proof}

\noindent
次の補題の証明方法は Gulliksen によるものらしい。

\begin{lemma}
\label{dpa-lemma-not-finite-pd}
$R' \to R$ を、その核が平方零で一つの元 $f$ によって生成される
Noether 環の全射とする。
$S = R[x_1, \ldots, x_n]/(f_1, \ldots, f_m)$.
$\sum r_j f_j = 0$ を $R[x_1, \ldots, x_n]$ における関係式とする。次を仮定する。
\begin{enumerate}
\item 各 $r_j$ の係数はすべて、$R$ における $f$ の零化イデアル $I$ に属する。
\item ある持ち上げ $r'_j, f'_j \in R'[x_1, \ldots, x_n]$ に対して
$\sum r'_j f'_j = gf$ であり、$g$ は $S/IS$ において冪零でない。
\end{enumerate}
このとき $S$ は $R$ 上有限 Tor 次元をもたない（すなわち $S$ は完全 $R$-代数でない）。
\end{lemma}

\begin{proof}
補題 \ref{dpa-lemma-tate-resolution} の Tate 分解 $R \to A \to S$ を選ぶ。
$\xi \in H_2(A/IA)$ および $\theta : A/IA \to A/IA$ を、補題
\ref{dpa-lemma-get-derivation} と \ref{dpa-lemma-compute-theta} で得られた元および導分とする。
$$
\theta^n(\gamma_n(\xi)) = g^n
$$
が $H_0(A/IA) = S/IS$ において成り立つことに注意する。
したがって $g$ が $S/IS$ において冪零でないならば、すべての $n > 0$ に対して
$\xi^n$ は $H_{2n}(A/IA)$ において非零である。
$H_{2n}(A/IA) = \text{Tor}^R_{2n}(S, R/I)$ なので結論が従う。
\end{proof}

\noindent
次の結果は \cite{Rodicio} に見いだされる。

\begin{lemma}
\label{dpa-lemma-injective}
$(A, \mathfrak m)$ を Noether 局所環とし、$I \subset J \subset A$ を真のイデアルとする。
$A/J$ が $A/I$ 上有限 Tor 次元をもつならば、
$I/\mathfrak m I \to J/\mathfrak m J$ は単射である。
\end{lemma}

\begin{proof}
$f \in I$ を、$I/\mathfrak m I$ における像が非零で、
$J/\mathfrak m J$ において零に写る元とする。
$\mathfrak m I \subset I' \subset I$ であり、$I/I'$ が $f$ の像によって生成されるような
イデアル $I'$ を選べる。$R = A/I$ および $R' = A/I'$ とおく。
ある $a_j \in A$ によって $J = (a_1, \ldots, a_m)$ とする。
すると、ある $b_j \in \mathfrak m$ に対して $f = \sum b_j a_j$ である。
$r_j, f_j \in R$ および $r'_j, f'_j \in R'$ を、それぞれ $b_j, a_j$ の像とする。
すると補題 \ref{dpa-lemma-not-finite-pd} の状況にある
（同補題のイデアル $I$ は $\mathfrak m_R$ に等しい）。これで補題は証明された。
\end{proof}

\begin{lemma}
\label{dpa-lemma-regular-sequence}
$(A, \mathfrak m)$ を Noether 局所環とし、$I \subset J \subset A$ を真のイデアルとする。
次を仮定する。
\begin{enumerate}
\item $A/J$ は $A/I$ 上有限 Tor 次元をもつ。
\item $J$ は正則列によって生成される。
\end{enumerate}
このとき $I$ は正則列によって生成され、$J/I$ も正則列によって生成される。
\end{lemma}

\begin{proof}
補題 \ref{dpa-lemma-injective} により、写像
$I/\mathfrak m I \to J/\mathfrak m J$ は単射である。したがって、$s \leq r$ と
$J$ の極小生成系 $f_1, \ldots, f_r$ を、$f_1, \ldots, f_s$ が $I$ に属し、
$I$ の極小生成系をなすように選べる。
More on Algebra, Lemmas
\ref{more-algebra-lemma-independence-of-generators} および
\ref{more-algebra-lemma-noetherian-finite-all-equivalent}
により、$J$ の任意の極小生成系は正則列であるから、補題が従う。
\end{proof}

\begin{lemma}
\label{dpa-lemma-perfect-map-ci}
$R \to S$ を Noether 局所環の局所環準同型とする。
$I \subset R$ および $J \subset S$ を $IS \subset J$ となるイデアルとする。
$R \to S$ が平坦で $S/\mathfrak m_R S$ が正則ならば、次の条件は同値である。
\begin{enumerate}
\item $J$ は正則列によって生成され、$S/J$ は $R/I$-加群として有限 Tor 次元をもつ。
\item $J$ は正則列によって生成され、
$\text{Tor}^{R/I}_p(S/J, R/\mathfrak m_R)$ が非零となる $p$ は有限個だけである。
\item $I$ は正則列によって生成され、$J/IS$ は $S/IS$ において正則列によって生成される。
\end{enumerate}
\end{lemma}

\begin{proof}
(3) が成り立つならば、$J$ は正則列によって生成される。例えば More on Algebra, Lemmas
\ref{more-algebra-lemma-join-koszul-regular-sequences} および
\ref{more-algebra-lemma-noetherian-finite-all-equivalent}
を参照せよ。さらに、(3) が成り立つならば、Koszul 複体が $S/IS$ 上の
$S/J$ の有限自由分解となるため、$S/J = (S/I)/(J/I)$ は $S/IS$ 上有限射影次元をもつ。
$R/I \to S/IS$ は平坦なので、More on Algebra, Lemma
\ref{more-algebra-lemma-flat-push-tor-amplitude} により、$S/J$ は $R/I$ 上有限 Tor 次元をもつ。
したがって (3) は (1) を含意する。

\medskip\noindent
(1) $\Rightarrow$ (2) は自明である。(2) を仮定する。More on Algebra, Lemma
\ref{more-algebra-lemma-perfect-over-regular-local-ring} により、$S/J$ は
$S/IS$ 上有限 Tor 次元をもつ。したがって補題
\ref{dpa-lemma-regular-sequence} を適用して、$IS$ および $J/IS$ が正則列によって生成されると分かる。
$f_1, \ldots, f_r \in I$ を $I$ の極小生成系とする。$R \to S$ は平坦なので、
$f_1, \ldots, f_r$ は $S$ において $IS$ の極小生成系をなす。
したがって、More on Algebra, Lemmas
\ref{more-algebra-lemma-independence-of-generators} および
\ref{more-algebra-lemma-noetherian-finite-all-equivalent}
により、元の列 $f_1, \ldots, f_r \in R$ の像は $S$ において正則列をなす。
ゆえに Algebra, Lemma \ref{algebra-lemma-flat-increases-depth} により、
$f_1, \ldots, f_r$ は $R$ における正則列である。
\end{proof}





\section{局所完全交叉環}
\label{dpa-section-lci}

\noindent
$(A, \mathfrak m)$ を Noether 完備局所環とする。Cohen 構造定理
（Algebra, Theorem \ref{algebra-theorem-cohen-structure-theorem} を参照）により、
$A$ はある正則 Noether 完備局所環 $R$ の商として書ける。
$R$ が正則局所環で、その核が正則列によって生成されるような全射
$R \to A$ が存在するとき、$A$ を{\it 完全交叉}という。
次の補題は、この概念が全射の選択によらないことを示す。

\begin{lemma}
\label{dpa-lemma-ci-well-defined}
$(A, \mathfrak m)$ を Noether 完備局所環とする。次の条件は同値である。
\begin{enumerate}
\item $R$ が正則局所環である任意の局所環の全射 $R \to A$ に対して、
$R \to A$ の核は正則列によって生成される。
\item $R$ が正則局所環であり、$R \to A$ の核が正則列によって生成されるような
局所環の全射 $R \to A$ が存在する。
\end{enumerate}
\end{lemma}

\begin{proof}
$k$ を $A$ の剰余体とする。$k$ の標数が $p > 0$ ならば、$\Lambda$ を
剰余体 $k$ をもつ Cohen 環とする（Algebra, Definition
\ref{algebra-definition-cohen-ring} および Lemma
\ref{algebra-lemma-cohen-rings-exist}）。$k$ の標数が $0$ ならば $\Lambda = k$ とおく。
任意の $n$ に対して $\Lambda[[x_1, \ldots, x_n]]$ は、$\mathfrak m$-進位相に関して
それぞれ $\mathbf{Z}$ または $\mathbf{Q}$ 上形式的に滑らかであることを思い出す。
More on Algebra, Lemma
\ref{more-algebra-lemma-power-series-ring-over-Cohen-fs}。
を参照せよ。Cohen 構造定理（Algebra, Theorem
\ref{algebra-theorem-cohen-structure-theorem}）における全射
$\Lambda[[x_1, \ldots, x_n]] \to A$ を一つ固定する。

\medskip\noindent
$R \to A$ を正則局所環 $R$ からの全射とし、$f_1, \ldots, f_r$ を
$\Ker(R \to A)$ の極小生成系とする。以下では、Noether 局所環のイデアルが
正則列によって生成されるための必要十分条件は、任意の極小生成系が正則列であることを、
断りなく用いる。Algebra, Lemmas
\ref{algebra-lemma-flat-increases-depth} および
\ref{algebra-lemma-completion-flat} により、$f_1, \ldots, f_r$ が $R$ において
正則列であるための必要十分条件は、完備化 $R^\wedge$ において正則列であることである。
さらに、
$$
R^\wedge/(f_1, \ldots, f_r)R^\wedge =
(R/(f_1, \ldots, f_n))^\wedge = A^\wedge = A
$$
である。実際、$A$ は $\mathfrak m_A$-進完備である（第一の等式は Algebra, Lemma
\ref{algebra-lemma-completion-tensor} による）。最後に、$R$ が正則なので
$R^\wedge$ も正則である（More on Algebra, Lemma
\ref{more-algebra-lemma-completion-regular}）。したがって $R$ は完備であると仮定してよい。

\medskip\noindent
$R$ が完備ならば、与えられた写像 $\Lambda[[x_1, \ldots, x_n]] \to A$ を
持ち上げる写像 $\Lambda[[x_1, \ldots, x_n]] \to R$ を選べる。
More on Algebra, Lemma \ref{more-algebra-lemma-lift-continuous} を参照せよ。
$R \to A$ の核の生成元に写る変数 $y_1, \ldots, y_m$ をさらに添加すれば、
$\Lambda[[x_1, \ldots, x_n, y_1, \ldots, y_m]] \to R$ が全射であると仮定できる
（詳細の一部は省略する）。すると可換図式
$$
\xymatrix{
\Lambda[[x_1, \ldots, x_n, y_1, \ldots, y_m]] \ar[r] \ar[d] & R \ar[d] \\
\Lambda[[x_1, \ldots, x_n]] \ar[r] & A
}
$$
を考えられる。Algebra, Lemma \ref{algebra-lemma-ci-well-defined} により、
$R \to A$ に対する条件は、固定した写像
$\Lambda[[x_1, \ldots, x_n]] \to A$ に対する条件と同値である。
これで補題の証明が完了する。
\end{proof}

\noindent
次の二つの補題は、上の定義に対する健全性検査である。

\begin{lemma}
\label{dpa-lemma-quotient-regular-ring-by-regular-sequence}
$R$ を正則環、$\mathfrak p \subset R$ を素イデアルとし、
$f_1, \ldots, f_r \in \mathfrak p$ を正則列とする。このとき
$$
A = (R/(f_1, \ldots, f_r))_\mathfrak p =
R_\mathfrak p/(f_1, \ldots, f_r)R_\mathfrak p
$$
の完備化は、上で定義した意味で完全交叉である。
\end{lemma}

\begin{proof}
$R_\mathfrak p$ は Noether 環で、その上の有限加群に対する完備化は完全なので
（Algebra, Lemma \ref{algebra-lemma-completion-tensor}）、$A$ の完備化は
$A^\wedge = R_\mathfrak p^\wedge/(f_1, \ldots, f_r)R_\mathfrak p^\wedge$
に等しい。Algebra, Lemmas \ref{algebra-lemma-completion-flat} および
\ref{algebra-lemma-flat-increases-depth} により、列 $f_1, \ldots, f_r$ の
$R_\mathfrak p$ における像は正則列である。さらに、More on Algebra, Lemma
\ref{more-algebra-lemma-completion-regular} により、$R_\mathfrak p^\wedge$ は正則局所環である。
したがって、完備局所環に対する完全交叉の定義から結果が従う。
\end{proof}

\noindent
次の補題は Algebra, Lemma \ref{algebra-lemma-lci} の類似である。

\begin{lemma}
\label{dpa-lemma-quotient-regular-ring}
$R$ を正則環、$\mathfrak p \subset R$ を素イデアル、
$I \subset \mathfrak p$ をイデアルとする。
$A = (R/I)_\mathfrak p = R_\mathfrak p/I_\mathfrak p$ とおく。次の条件は同値である。
\begin{enumerate}
\item $A$ の完備化は上の意味で完全交叉である。
\item $I_\mathfrak p \subset R_\mathfrak p$ は正則列によって生成される。
\item 加群 $(I/I^2)_\mathfrak p$ は
$\dim(R_\mathfrak p) - \dim(A)$ 個の元によって生成できる。
\item ここにさらに追加する。
\end{enumerate}
\end{lemma}

\begin{proof}
$R$ を $\mathfrak p$ における局所化に置き換えてよく、実際そうする。
すると $\mathfrak p = \mathfrak m$ は $R$ の極大イデアルで、$A = R/I$ である。
$f_1, \ldots, f_r \in I$ を極小生成系とする。Noether 環 $R_\mathfrak p$ 上の
有限加群に対する完備化は完全なので（Algebra, Lemma
\ref{algebra-lemma-completion-tensor}）、$A$ の完備化は
$A^\wedge = R^\wedge/(f_1, \ldots, f_r)R^\wedge$
に等しい。

\medskip\noindent
(1) が成り立つならば、仮定により列 $f_1, \ldots, f_r$ の $R^\wedge$ における像は
正則列である。したがって Algebra, Lemmas \ref{algebra-lemma-completion-flat} および
\ref{algebra-lemma-flat-increases-depth} により、これは $R$ においても正則列である。
ゆえに (1) は (2) を含意する。

\medskip\noindent
(3) を仮定する。$c = \dim(R) - \dim(A)$ とおき、その像が $I/I^2$ を生成する
$f_1, \ldots, f_c \in I$ を選ぶ。Nakayama の補題
（Algebra, Lemma \ref{algebra-lemma-NAK}）により、
$I = (f_1, \ldots, f_c)$ である。$R$ は正則、したがって Cohen--Macaulay なので
（Algebra, Proposition \ref{algebra-proposition-CM-module}）、Algebra, Proposition
\ref{algebra-proposition-CM-module} により $f_1, \ldots, f_c$ は正則列である。
したがって (3) は (2) を含意する。最後に、補題
\ref{dpa-lemma-quotient-regular-ring-by-regular-sequence} により (2) は (1) を含意する。
\end{proof}

\noindent
次の結果は Avramov による。\cite{Avramov} を参照せよ。

\begin{proposition}
\label{dpa-proposition-avramov}
$A \to B$ を Noether 局所環の平坦な局所準同型とする。次の条件は同値である。
\begin{enumerate}
\item $B^\wedge$ は完全交叉である。
\item $A^\wedge$ および $(B/\mathfrak m_A B)^\wedge$ は完全交叉である。
\end{enumerate}
\end{proposition}

\begin{proof}
図式
$$
\xymatrix{
B \ar[r] & B^\wedge \\
A \ar[u] \ar[r] & A^\wedge \ar[u]
}
$$
を考える。水平写像は忠実平坦なので
（Algebra, Lemma \ref{algebra-lemma-completion-faithfully-flat}）、
右の垂直矢印は平坦である（例えば Algebra, Lemma
\ref{algebra-lemma-criterion-flatness-fibre-Noetherian} による）。さらに、
$(B/\mathfrak m_A B)^\wedge = B^\wedge/\mathfrak m_{A^\wedge} B^\wedge$
が Algebra, Lemma \ref{algebra-lemma-completion-tensor} により成り立つ。
したがって $A$ と $B$ は完備 Noether 局所環であると仮定してよい。

\medskip\noindent
$A$ と $B$ を完備 Noether 局所環と仮定する。More on Algebra, Lemma
\ref{more-algebra-lemma-embed-map-Noetherian-complete-local-rings}
のような図式
$$
\xymatrix{
S \ar[r] & B \\
R \ar[u] \ar[r] & A \ar[u]
}
$$
を選ぶ。$I = \Ker(R \to A)$ および $J = \Ker(S \to B)$ とおく。
$R/I = A \to B = S/J$ は平坦なので、写像
$J/IS \otimes_R R/\mathfrak m_R \to J/J \cap \mathfrak m_R S$
は同型であることに注意する。したがって $J/IS$ の極小生成系は
$\Ker(S/\mathfrak m_R S \to B/\mathfrak m_A B)$ の極小生成系に写る。
最後に、$S/\mathfrak m_R S$ は正則局所環である。

\medskip\noindent
(1) が成り立つ、すなわち $J$ が正則列によって生成されると仮定する。
$A = R/I \to B = S/J$ は平坦なので補題 \ref{dpa-lemma-perfect-map-ci} が適用でき、
$I$ と $J/IS$ は正則列によって生成されると分かる。
$\dim(B) = \dim(A) + \dim(B/\mathfrak m_A B)$ および
$\dim(S/IS) = \dim(A) + \dim(S/\mathfrak m_R S)$ である
（Algebra, Lemma \ref{algebra-lemma-dimension-base-fibre-equals-total}）。
したがって $J/IS$ は
$$
\dim(S/IS) - \dim(S/J) = \dim(S/\mathfrak m_R S) - \dim(B/\mathfrak m_A B)
$$
個の元によって生成される（Algebra, Lemma \ref{algebra-lemma-one-equation}）。
したがって $\Ker(S/\mathfrak m_R S \to B/\mathfrak m_A B)$ も同じ個数の元によって
生成される（上を参照）。ゆえにこの核は正則列によって生成される。
例えば補題 \ref{dpa-lemma-quotient-regular-ring} を参照せよ。こうして (2) が成り立つと分かる。

\medskip\noindent
(2) が成り立つならば、$I$ および $J/J \cap \mathfrak m_R S$ は正則列によって生成される。
これらの生成元を持ち上げ（上を参照）、$R/I \to S/IS$ の平坦性と
Grothendieck の補題（Algebra, Lemma
\ref{algebra-lemma-grothendieck-regular-sequence}）を用いると、$J/IS$ は
$S/IS$ において正則列によって生成されると分かる。すると補題
\ref{dpa-lemma-perfect-map-ci} により $J$ は正則列によって生成されるので、(1) が成り立つ。
\end{proof}

\begin{definition}
\label{dpa-definition-lci}
$A$ を Noether 環とする。
\begin{enumerate}
\item $A$ が局所環で、その完備化が上の意味で完全交叉であるとき、
$A$ を{\it 完全交叉}という。
\item 一般に、$A$ のすべての局所環が完全交叉であるとき、
$A$ を{\it 局所完全交叉}という。
\end{enumerate}
\end{definition}

\noindent
以下で、これが Algebra, Definitions \ref{algebra-definition-lci-field} および
\ref{algebra-definition-lci-local-ring} で導入した用語と矛盾しないことを確認する。
しかしまず、$A$ が Noether 局所完全交叉ならば $A$ は局所完全交叉、
すなわちすべての局所環が完全交叉であることを示し、この定義が「意味をなす」ことを確かめる。

\begin{lemma}
\label{dpa-lemma-ci-good}
$(A, \mathfrak m)$ を Noether 局所環とし、$\mathfrak p \subset A$ を素イデアルとする。
$A$ が完全交叉ならば、$A_\mathfrak p$ も完全交叉である。
\end{lemma}

\begin{proof}
$\mathfrak p$ の上にある $A^\wedge$ の素イデアル $\mathfrak q$ を選ぶ。
これは可能である。実際、Algebra, Lemma
\ref{algebra-lemma-completion-faithfully-flat} により $A \to A^\wedge$ は忠実平坦である。
すると $A_\mathfrak p \to (A^\wedge)_\mathfrak q$ は平坦な局所環準同型である。
したがって命題 \ref{dpa-proposition-avramov} により、$A_\mathfrak p$ が完全交叉であるための
必要十分条件は $(A^\wedge)_\mathfrak q$ が完全交叉であることである。
ゆえに $A$ が完備の場合に補題を証明すれば十分である（これが証明の要点である）。

\medskip\noindent
$A$ が完備であると仮定する。定義により、正則局所環 $R$ の正則列
$f_1, \ldots, f_r$ を用いて $A = R/(f_1, \ldots, f_r)$ と書ける。
$\mathfrak q \subset R$ を $\mathfrak p$ に対応する素イデアルとする。
$f_1, \ldots, f_r \in \mathfrak q$ であり、
$A_\mathfrak p = R_\mathfrak q/(f_1, \ldots, f_r)R_\mathfrak q$ であることに注意する。
したがって補題 \ref{dpa-lemma-quotient-regular-ring-by-regular-sequence} により、
$A_\mathfrak p$ は完全交叉である。
\end{proof}

\begin{lemma}
\label{dpa-lemma-check-lci-at-maximal-ideals}
$A$ を Noether 環とする。$A$ が局所完全交叉であるための必要十分条件は、
$A$ のすべての極大イデアル $\mathfrak m$ に対して $A_\mathfrak m$ が完全交叉であることである。
\end{lemma}

\begin{proof}
補題 \ref{dpa-lemma-ci-good} と定義から直ちに従う。
\end{proof}

\begin{lemma}
\label{dpa-lemma-check-lci-agrees}
$S$ を体 $k$ 上の有限型代数とする。
\begin{enumerate}
\item 素イデアル $\mathfrak q \subset S$ に対して、局所環 $S_\mathfrak q$ が
Algebra, Definition \ref{algebra-definition-lci-local-ring} の意味で完全交叉であるための
必要十分条件は、定義 \ref{dpa-definition-lci} の意味で完全交叉であることである。
\item $S$ が Algebra, Definition \ref{algebra-definition-lci-field} の意味で
局所完全交叉であるための必要十分条件は、定義 \ref{dpa-definition-lci} の意味で
局所完全交叉であることである。
\end{enumerate}
\end{lemma}

\begin{proof}
(1) を証明する。$k[x_1, \ldots, x_n] \to S$ を全射とする。
$\mathfrak p \subset k[x_1, \ldots, x_n]$ を $\mathfrak q$ に対応する素イデアルとし、
$I \subset k[x_1, \ldots, x_n]$ をこの全射の核とする。
$k[x_1, \ldots, x_n]_\mathfrak p \to S_\mathfrak q$ は核 $I_\mathfrak p$ をもつ全射である。
Algebra, Proposition \ref{algebra-proposition-finite-gl-dim-polynomial-ring} により、
$k[x_1, \ldots, x_n]$ は正則環である。したがって補題
\ref{dpa-lemma-quotient-regular-ring} と Algebra, Lemma
\ref{algebra-lemma-lci-local} を組み合わせると、(1) の二つの概念の同値性が従う。

\medskip\noindent
(1) を証明したので、(2) の同値性は定義と Algebra, Lemma
\ref{algebra-lemma-lci-global} から従う。
\end{proof}

\begin{lemma}
\label{dpa-lemma-avramov}
$A \to B$ を Noether 局所環の平坦な局所準同型とする。次の条件は同値である。
\begin{enumerate}
\item $B$ は完全交叉である。
\item $A$ および $B/\mathfrak m_A B$ は完全交叉である。
\end{enumerate}
\end{lemma}

\begin{proof}
定義が意味をなすことが分かったので、これは（非自明な）命題
\ref{dpa-proposition-avramov} の自明な言い換えである。
\end{proof}









\section{局所完全交叉準同型}
\label{dpa-section-lci-homomorphisms}

\noindent
$A \to B$ を Noether 完備局所環の局所準同型とする。Cohen 構造定理の帰結として、
Noether 完備局所環の可換図式
$$
\xymatrix{
S \ar[r] & B \\
& A \ar[lu] \ar[u]
}
$$
で、$S \to B$ が全射、$A \to S$ が平坦、かつ
$S/\mathfrak m_A S$ が正則局所環となるものが見つかる。
これは More on Algebra, Lemma
\ref{more-algebra-lemma-embed-map-Noetherian-complete-local-rings} から従う。
一時的に、$S$ が Noether 局所環、$A \to S \to B$ が局所環準同型、
$S \to B$ が全射、$A \to S$ が平坦、かつ $S/\mathfrak m_A S$ が正則であるとき、
$A \to S \to B$ を $A \to B$ の{\it 良好な分解}という。
$S \to B$ の核が正則列によって生成されるような良好な分解
$A \to S \to B$ が存在するとき、$A \to B$ を{\it 完全交叉準同型}という。
次の補題は、この概念が図式の選択によらないことを示す。

\begin{lemma}
\label{dpa-lemma-ci-map-well-defined}
$A \to B$ を Noether 完備局所環の局所準同型とする。次の条件は同値である。
\begin{enumerate}
\item ある良好な分解 $A \to S \to B$ に対して、$S \to B$ の核は正則列によって生成される。
\item すべての良好な分解 $A \to S \to B$ に対して、$S \to B$ の核は正則列によって生成される。
\end{enumerate}
\end{lemma}

\begin{proof}
$A \to S \to B$ を良好な分解とする。$B$ は完備なので、$S^\wedge$ を $S$ の完備化として
分解 $A \to S^\wedge \to B$ を得る。これも良好な分解であることに注意する。
環準同型 $S \to S^\wedge$ は平坦なので
（Algebra, Lemma \ref{algebra-lemma-completion-flat}）、$A \to S^\wedge$ は平坦である。
$S/\mathfrak m_A S$ が正則なので、環
$S^\wedge/\mathfrak m_A S^\wedge = (S/\mathfrak m_A S)^\wedge$ も正則である
（More on Algebra, Lemma \ref{more-algebra-lemma-completion-regular}）。
$f_1, \ldots, f_r$ を $\Ker(S \to B)$ の極小生成系とする。
以下では、Noether 局所環のイデアルが正則列によって生成されるための必要十分条件は、
任意の極小生成系が正則列であることを断りなく用いる。Algebra, Lemma
\ref{algebra-lemma-flat-increases-depth} により、$f_1, \ldots, f_r$ が $S$ において
正則列であるための必要十分条件は、完備化 $S^\wedge$ において正則列であることである。
さらに、
$$
S^\wedge/(f_1, \ldots, f_r)R^\wedge =
(S/(f_1, \ldots, f_n))^\wedge = B^\wedge = B
$$
である。実際、$B$ は $\mathfrak m_B$-進完備である（第一の等式は Algebra, Lemma
\ref{algebra-lemma-completion-tensor} による）。したがって $S \to B$ の核が正則列によって
生成されるための必要十分条件は、$S^\wedge \to B$ の核が正則列によって生成されることである。
ゆえに $S$ が完備である良好な分解だけを考えれば十分である。

\medskip\noindent
$S$ と $S'$ が完備である二つの分解 $A \to S \to B$ および
$A \to S' \to B$ があると仮定する。More on Algebra, Lemma
\ref{more-algebra-lemma-dominate-two-surjections} により、環
$S \times_B S'$ は Noether 完備局所環である。したがって More on Algebra, Lemma
\ref{more-algebra-lemma-embed-map-Noetherian-complete-local-rings} を用いて、
$S''$ が完備である良好な分解 $A \to S'' \to S \times_B S'$ を選べる。
ゆえに、$A \to S' \to S \to B$ が比較可能な良好な分解ならば、
$\Ker(S \to B)$ が正則列によって生成されるための必要十分条件は、
$\Ker(S' \to B)$ が正則列によって生成されることである、と示せば十分である。

\medskip\noindent
$A \to S' \to S \to B$ を比較可能な良好な分解とする。まず、
$S'/\mathfrak m_R S' \to S/\mathfrak m_R S$ は正則局所環の全射なので、その核は正則列
$\overline{x}_1, \ldots, \overline{x}_c \in
\mathfrak m_{S'}/\mathfrak m_R S'$
によって生成される。この列は正則局所環 $S'/\mathfrak m_R S'$ の
正則パラメータ系に延長できる。Algebra, Lemma
\ref{algebra-lemma-regular-quotient-regular} を参照せよ。
$I = \Ker(S' \to S)$ とおく。$S$ は $R$ 上平坦なので、
$$
I/\mathfrak m_R I =
\Ker(S'/\mathfrak m_R S' \to S/\mathfrak m_R S) =
(\overline{x}_1, \ldots, \overline{x}_c).
$$
持ち上げ $x_1, \ldots, x_c \in I$ を選ぶ。$S'$ は $R$ 上平坦なので、
これらの持ち上げは $I$ を生成する正則列をなす。
Algebra, Lemma \ref{algebra-lemma-grothendieck-regular-sequence} を参照せよ。

\medskip\noindent
したがって、$\Ker(S \to B)$ も正則列によって生成されるならば、
$\Ker(S' \to B)$ も同様である。More on Algebra, Lemmas
\ref{more-algebra-lemma-join-koszul-regular-sequences} および
\ref{more-algebra-lemma-noetherian-finite-all-equivalent} を参照せよ。

\medskip\noindent
逆に、$J = \Ker(S' \to B)$ が正則列によって生成されると仮定する。
$I$ の生成元 $x_1, \ldots, x_c$ は $\mathfrak m_{S'}/\mathfrak m_{S'}^2$ において
線形独立な元に写るので、写像
$I/\mathfrak m_{S'}I \to J/\mathfrak m_{S'}J$ は単射である。
したがって $J$ は極小生成系 $x_1, \ldots, x_c, y_1, \ldots, y_d$ をもつ。
この列は正則列であり、したがって $y_1, \ldots, y_d$ の $S$ における像は
$\Ker(S \to B)$ を生成する正則列をなす。これで補題の証明が完了する。
\end{proof}

\noindent
次の命題では、Tor の消滅に関する条件は、特に $B$ が $A$ 上有限 Tor 次元をもつ場合、
したがって特に $B$ が $A$ 上平坦な場合に成り立つことに注意する。

\begin{proposition}
\label{dpa-proposition-avramov-map}
$A \to B$ を Noether 局所環の局所準同型とする。次の条件は同値である。
\begin{enumerate}
\item $B$ は完全交叉で、$\text{Tor}^A_p(B, A/\mathfrak m_A)$ が非零となる $p$ は
有限個だけである。
\item $A$ は完全交叉で、$A^\wedge \to B^\wedge$ は上で定義した意味で
完全交叉準同型である。
\end{enumerate}
\end{proposition}

\begin{proof}
$F_\bullet \to A/\mathfrak m_A$ を有限自由 $A$-加群による分解とする。
$\text{Tor}^A_p(B, A/\mathfrak m_A)$ は複体 $F_\bullet \otimes_A B$ の
第 $p$ ホモロジーであることに注意する。完備化
$F_\bullet^\wedge = F_\bullet \otimes_A A^\wedge$ とおく。
すると $F_\bullet^\wedge$ は、有限自由 $A^\wedge$-加群による
$A^\wedge/\mathfrak m_{A^\wedge}$ の分解である。実際、$A \to A^\wedge$ は平坦で、
有限加群上の完備化は完全である。Algebra, Lemmas
\ref{algebra-lemma-completion-tensor} および
\ref{algebra-lemma-completion-flat} を参照せよ。したがって
$$
F_\bullet^\wedge \otimes_{A^\wedge} B^\wedge =
F_\bullet \otimes_A B \otimes_B B^\wedge
$$
$B \to B^\wedge$ の平坦性により、
$$
\text{Tor}^{A^\wedge}_p(B^\wedge, A^\wedge/\mathfrak m_{A^\wedge}) =
\text{Tor}^A_p(B, A/\mathfrak m_A) \otimes_B B^\wedge
$$
を得る。こうして、局所環準同型 $A \to B$ に対する (1) の条件は、
局所環準同型 $A^\wedge \to B^\wedge$ に対する同じ条件と同値であると分かる。
したがって $A$ と $B$ は完備 Noether 局所環であると仮定してよい
（いずれにせよ他の条件は完備化を用いて定式化されている）。

\medskip\noindent
$A$ と $B$ を完備 Noether 局所環と仮定する。More on Algebra, Lemma
\ref{more-algebra-lemma-embed-map-Noetherian-complete-local-rings}
のような図式
$$
\xymatrix{
S \ar[r] & B \\
R \ar[u] \ar[r] & A \ar[u]
}
$$
を選ぶ。$I = \Ker(R \to A)$ および $J = \Ker(S \to B)$ とおく。
この命題は補題 \ref{dpa-lemma-perfect-map-ci} から従う。
\end{proof}

\begin{remark}
\label{dpa-remark-no-good-ci-map}
一般の Noether 環の間の写像に対して「局所完全交叉準同型」の良い概念を
定義することは困難に思われる。その理由は、Noether 局所環 $A$ に対して
$A \to A^\wedge$ のファイバーが局所完全交叉環ではないからである。
したがって、$A \to B$ が Noether 局所環の局所準同型で、完備化の写像
$A^\wedge \to B^\wedge$ が上で定義した意味の完全交叉準同型であっても、一般には
$(A_\mathfrak p)^\wedge \to (B_\mathfrak q)^\wedge$ は上で定義した意味の
完全交叉準同型では{\bf ない}。優秀 Noether 環だけを扱えば解決できる。
より一般には、形式的ファイバーが完全交叉である Noether 環を扱うことができる。
\cite{Rodicio-ci} を参照せよ。この理論は Dualizing Complexes, Section
\ref{dualizing-section-formal-fibres} で展開する。
\end{remark}

\noindent
本節の最後に、上で定義した概念と More on Algebra, Section
\ref{dpa-section-lci} で導入した概念とを比較する。

\begin{lemma}
\label{dpa-lemma-well-defined-if-you-can-find-good-factorization}
Noether 局所環の可換図式
$$
\xymatrix{
S \ar[r] & B \\
& A \ar[lu] \ar[u]
}
$$
で、$S \to B$ が全射、$A \to S$ が平坦、かつ
$S/\mathfrak m_A S$ が正則局所環であるものを考える。次の条件は同値である。
\begin{enumerate}
\item $\Ker(S \to B)$ は正則列によって生成される。
\item $A^\wedge \to B^\wedge$ は上で定義した完全交叉準同型である。
\end{enumerate}
\end{lemma}

\begin{proof}
省略する。ヒント：証明は補題 \ref{dpa-lemma-ci-map-well-defined} の証明の
第一段落の議論と同じである。
\end{proof}

\begin{lemma}
\label{dpa-lemma-finite-type-lci-map}
$A$ を Noether 環とし、$A \to B$ を有限型環準同型とする。次の条件は同値である。
\begin{enumerate}
\item $A \to B$ は More on Algebra, Definition
\ref{more-algebra-definition-local-complete-intersection} の意味で局所完全交叉である。
\item すべての素イデアル $\mathfrak q \subset B$ に対して、
$\mathfrak p = A \cap \mathfrak q$ とおくと、環準同型
$(A_\mathfrak p)^\wedge \to (B_\mathfrak q)^\wedge$ は上で定義した意味で
完全交叉準同型である。
\end{enumerate}
\end{lemma}

\begin{proof}
全射 $R = A[x_1, \ldots, x_n] \to B$ を選ぶ。$A \to R$ は平坦で、
ファイバーは正則であることに注意する。$I$ を $R \to B$ の核とする。
(2) を仮定する。補題 \ref{dpa-lemma-well-defined-if-you-can-find-good-factorization} および
Algebra, Lemma \ref{algebra-lemma-regular-sequence-in-neighbourhood} により、
$I$ は局所的に正則列によって生成される。言い換えれば (1) が成り立つ。
逆に (1) を仮定する。$R$ と $B$ を局所化した後、$I$ は Koszul 正則列によって
生成されると仮定できる。More on Algebra, Lemma
\ref{more-algebra-lemma-noetherian-finite-all-equivalent} により、$I$ は局所的に
正則列によって生成される。したがって補題
\ref{dpa-lemma-well-defined-if-you-can-find-good-factorization} により (2) が成り立つ。
詳細の一部は省略する。
\end{proof}

\begin{lemma}
\label{dpa-lemma-avramov-map-finite-type}
$A$ を Noether 環とする。$A \to B$ を有限型環準同型で、
$\Spec(B) \to \Spec(A)$ の像が $\Spec(A)$ のすべての閉点を含むものとする。
次の条件は同値である。
\begin{enumerate}
\item $B$ は完全交叉で、$A \to B$ は有限 Tor 次元をもつ。
\item $A$ は完全交叉で、$A \to B$ は More on Algebra, Definition
\ref{more-algebra-definition-local-complete-intersection} の意味で局所完全交叉である。
\end{enumerate}
\end{lemma}

\begin{proof}
これは命題 \ref{dpa-proposition-avramov-map} を補題
\ref{dpa-lemma-finite-type-lci-map} により言い換えたものである。詳細は省略する。
\end{proof}








\section{滑らかな環準同型と対角射}
\label{dpa-section-smooth-diagonal-perfect}

\noindent
本節では上の内容を用いて、滑らかな環準同型を、対角射が完全であるものとして特徴付ける。

\begin{lemma}
\label{dpa-lemma-local-perfect-diagonal}
$A \to B$ を Noether 局所環の局所準同型で、$B$ が $A$ 上平坦かつ本質的有限型であるものとする。
$$
B \otimes_A B \longrightarrow B
$$
が完全な環準同型、すなわち $B$ が $B \otimes_A B$ 上有限 Tor 次元をもつならば、
$B$ は滑らかな $A$-代数の局所化である。
\end{lemma}

\begin{proof}
$B$ は $A$ 上本質的有限型なので、$B \otimes_A B$ も同様であり、特に
$B \otimes_A B$ は Noether 環である。したがって、$B \otimes_A B$ の商 $B$ は
$B \otimes_A B$ 上擬コヒーレントである
（More on Algebra, Lemma \ref{more-algebra-lemma-Noetherian-pseudo-coherent}）。
これにより、環準同型の完全性（More on Algebra, Definition
\ref{more-algebra-definition-pseudo-coherent-perfect}）が有限 Tor 次元という条件と一致する理由が分かる。

\medskip\noindent
$B = R/K$ と書ける。ここで $R$ は $A[x_1, \ldots, x_n]$ のある素イデアルにおける
局所化で、$K \subset R$ はイデアルである。$\mathfrak m \subset R \otimes_A R$ を、全射
$R \otimes_A R \to B \otimes_A B \to B$ による $B$ の極大イデアルの逆像である
極大イデアルとする。すると全射
$$
(R \otimes_A R)_\mathfrak m \to (B \otimes_A B)_\mathfrak m \to B
$$
があり、したがって補題 \ref{dpa-lemma-injective} のようなイデアル
$I \subset J \subset (R \otimes_A R)_\mathfrak m$ が得られる。
$I/\mathfrak m I \to J/\mathfrak m J$ は単射であると結論される。

\medskip\noindent
$K = (f_1, \ldots, f_r)$ と書き、$r$ は極小とする。$f_i \in R$ は
$A[x_1, \ldots, x_n]$ の元の像であると仮定してよく、実際そうする。その元も $f_i$ と表す。
$I$ は $f_1 \otimes 1, \ldots, f_r \otimes 1$ および
$1 \otimes f_1, \ldots, 1 \otimes f_r$ によって生成されることに注意する。
これは $I$ の極小生成系であると主張する。実際、$\kappa$ を $R$、$B$、
$(R \otimes_A R)_\mathfrak m$、および $(B \otimes_A B)_\mathfrak m$ の共通の剰余体とすると、写像
$R \otimes_A R \to R \otimes_A \kappa \oplus \kappa \otimes_A R$
があり、これは $(R \otimes_A R)_\mathfrak m$ を経由する。
$B$ は $A$ 上平坦で、短完全列 $0 \to K \to R \to B \to 0$ があるので、
$K \otimes_A \kappa \subset R \otimes_A \kappa$ である。
Algebra, Lemma \ref{algebra-lemma-flat-tor-zero} を参照せよ。したがって写像
$(R \otimes_A R)_\mathfrak m \to R \otimes_A \kappa \oplus \kappa \otimes_A R$
を $I$ に制限すると、写像
$$
I \to K \otimes_A \kappa \oplus \kappa \otimes_A K \to
K \otimes_B \kappa \oplus \kappa \otimes_B K.
$$
を得る。元
$f_1 \otimes 1, \ldots, f_r \otimes 1, 1 \otimes f_1, \ldots, 1 \otimes f_r$
はこの写像の終域の基に写る。実際、Nakayama の補題
（Algebra, Lemma \ref{algebra-lemma-NAK}）により、$f_1, \ldots, f_r$ は
$K \otimes_B \kappa$ の基に写る。これで主張が証明された。

\medskip\noindent
イデアル $J$ は $f_1 \otimes 1, \ldots, f_r \otimes 1$ および元
$x_1 \otimes 1 - 1 \otimes x_1, \ldots, x_n \otimes 1 - 1 \otimes x_n$
によって生成される（証明には、これらの元がイデアル $J$ に属することだけ分かれば十分である）。
ここで
$$
f_i \otimes 1 - 1 \otimes f_i =
\sum g_{ij} (x_j \otimes 1 - 1 \otimes x_j)
$$
と書ける。ただし $g_{ij} \in (R \otimes_A R)_\mathfrak m$ である。
これは $A[x_1, \ldots, x_n]$ の元に関する一般的事実であり、証明は省略する。
$a_{ij} \in \kappa$ を $g_{ij}$ の像とする。別の計算により、$a_{ij}$ は
$\partial f_i/\partial x_j$ の $\kappa$ における像であることが分かる。
$I/\mathfrak m I \to J/\mathfrak m J$ の単射性と上の注意から、行列
$(a_{ij})$ は最大階数 $r$ をもたなければならない。
$$
C = A[x_1, \ldots, x_n]/(f_1, \ldots, f_r)
$$
とおき、素朴な余接複体
$\NL_{C/A} \cong (C^{\oplus r} \to C^{\oplus n})$ を考える。ここで写像は偏微分の行列で与えられる。
したがって $\NL_{C/A} \otimes_A B$ は、次数 $0$ に置かれた階数 $n - r$ の
自由 $B$-加群と同型である。ゆえに $B$ において単元に写るある $g \in C$ に対して、
$C_g$ は $A$ 上滑らかである。Algebra, Lemma
\ref{algebra-lemma-smooth-at-point} を参照せよ。これで証明が完了する。
\end{proof}

\begin{lemma}
\label{dpa-lemma-perfect-diagonal}
$A \to B$ を Noether 環の平坦な有限型環準同型とする。
$$
B \otimes_A B \longrightarrow B
$$
が完全な環準同型、すなわち $B$ が $B \otimes_A B$ 上有限 Tor 次元をもつならば、
$B$ は滑らかな $A$-代数である。
\end{lemma}

\begin{proof}
これは補題 \ref{dpa-lemma-local-perfect-diagonal} と滑らかな環準同型に関する一般的事実から従う。
Algebra, Lemmas \ref{algebra-lemma-smooth-at-point} および
\ref{algebra-lemma-locally-smooth} を参照せよ。別の方法として、読者は補題
\ref{dpa-lemma-local-perfect-diagonal} の証明を少し修正してこの補題を証明してもよい。
\end{proof}






\section{余法加群の自由性}
\label{dpa-section-freeness-conormal}

\noindent
Tate 分解とその上の導分を用いて、本節の結果の（より強い）版を証明できる。
\cite{Iyengar} を参照せよ。より初等的な二つの参考文献は
\cite{Vasconcelos} および \cite{Ferrand-lci} である。

\begin{lemma}
\label{dpa-lemma-free-summand-in-ideal-finite-proj-dim}
\begin{reference}
\cite{Vasconcelos}
\end{reference}
$R$ を Noether 局所環とする。$I \subset R$ を $R$ 上有限射影次元をもつイデアルとする。
$F \subset I/I^2$ が $R/I$ と同型な直和因子ならば、$I/I^2$ における像が $F$ を
生成するような非零因子 $x \in I$ が存在する。
\end{lemma}

\begin{proof}
仮定により、有限自由分解
$$
0 \to R^{\oplus n_e} \to R^{\oplus n_{e-1}} \to \ldots \to
R^{\oplus n_1} \to R \to R/I \to 0
$$
を選べる。このとき $\varphi_1 : R^{\oplus n_1} \to R$ は階数 $1$ であり、
Algebra, Proposition \ref{algebra-proposition-what-exact} により、$I$ は非零因子 $y$ を含む。
$\mathfrak p_1, \ldots, \mathfrak p_n$ を $R$ の随伴素イデアルとする。
Algebra, Lemma \ref{algebra-lemma-finite-ass} を参照せよ。
$J/I^2 = F$ となるイデアル $I^2 \subset J \subset I$ を選ぶ。
$y^2 \in J$ かつ $y^2 \not\in \mathfrak p_i$ なので、すべての $i$ に対して
$J \not\subset \mathfrak p_i$ である。Algebra, Lemma
\ref{algebra-lemma-ass-zero-divisors} を参照せよ。
Nakayama の補題（Algebra, Lemma \ref{algebra-lemma-NAK}）により、
$J \not\subset \mathfrak m J + I^2$ である。Algebra, Lemma
\ref{algebra-lemma-silly} により、$x \in J$ を、
$x \not\in \mathfrak m J + I^2$ かつ $i = 1, \ldots, n$ に対して
$x \not\in \mathfrak p_i$ となるように選べる。すると $x$ は非零因子であり、
$I/I^2$ における $x$ の像は、Nakayama の補題により直和因子
$J/I^2 \cong R/I$ を生成する。
\end{proof}

\begin{lemma}
\label{dpa-lemma-vasconcelos}
\begin{reference}
\cite[Theorem 1.1]{Vasconcelos} の局所版
\end{reference}
$R$ を Noether 局所環とする。$I \subset R$ を $R$ 上有限射影次元をもつイデアルとする。
$F \subset I/I^2$ が階数 $r$ の自由直和因子ならば、
$x_1 \bmod I^2, \ldots, x_r \bmod I^2$ が $F$ を生成するような正則列
$x_1, \ldots, x_r \in I$ が存在する。
\end{lemma}

\begin{proof}
$r = 0$ ならば示すことはない。$r > 0$ と仮定する。補題
\ref{dpa-lemma-free-summand-in-ideal-finite-proj-dim} により、$x \in I$ を、
$x$ が非零因子で、$x \bmod I^2$ が $F$ の中の $R/I$ と同型な直和因子を
生成するように選べる。環 $R' = R/(x)$ とイデアル $I' = I/(x)$ を考える。
もちろん $R'/I' = R/I$ である。短完全列
$$
0 \to R/I \xrightarrow{x} I/xI \to I' \to 0
$$
は分裂する。実際、写像 $I/xI \to I/I^2$ は $xR/xI$ を直和因子に写す。
ここで $I/xI = I \otimes_R^\mathbf{L} R'$ は $R'$ 上有限射影次元をもつ。
More on Algebra, Lemmas \ref{more-algebra-lemma-perfect-module} および
\ref{more-algebra-lemma-pull-perfect} を参照せよ。
したがって直和因子 $I'$ は $R'$ 上有限射影次元をもつ。一方、短完全列
$0 \to xR/xI \to I/I^2 \to I'/(I')^2 \to 0$ があるので、
$I'/(I')^2$ は階数 $r - 1$ の自由直和因子 $F' = F/(R/I \cdot x)$ をもつ。
$r$ に関する帰納法により、その合同類が $F'$ を自由生成する正則列
$x'_2, \ldots, x'_r \in I'$ を選べる。$x_1 = x$ とし、$x_2, \ldots, x_r$ を
$x'_1, \ldots, x'_r$ の $R$ における任意の持ち上げとすれば、補題が成り立つと分かる。
\end{proof}

\begin{proposition}
\label{dpa-proposition-regular-ideal}
\begin{reference}
\cite[Corollary 1]{Vasconcelos} の変形版。さらに
\cite{Iyengar} および \cite{Ferrand-lci} を参照せよ。
\end{reference}
$R$ を Noether 環とする。$I \subset R$ を有限射影次元をもち、$I/I^2$ が
$R/I$ 上有限局所自由であるイデアルとする。このとき $I$ は正則イデアルである
（More on Algebra, Definition \ref{more-algebra-definition-regular-ideal}）。
\end{proposition}

\begin{proof}
Algebra, Lemma \ref{algebra-lemma-regular-sequence-in-neighbourhood} により、
すべての $\mathfrak p \supset I$ に対して $I_\mathfrak p \subset R_\mathfrak p$ が
正則列によって生成されることを示せば十分である。したがって $R$ は局所環と仮定してよい。
$I/I^2$ の階数が $r$ ならば、補題 \ref{dpa-lemma-vasconcelos} により、
$I/I^2$ を生成する正則列 $x_1, \ldots, x_r \in I$ が見つかる。
Nakayama の補題（Algebra, Lemma \ref{algebra-lemma-NAK}）により、
$I$ は $x_1, \ldots, x_r$ によって生成される。
\end{proof}

\noindent
任意の環の局所完全交叉準同型 $A \to B$ に対して、素朴な余接複体 $\NL_{B/A}$ は
Tor 振幅が $[-1, 0]$ にある完全複体である。More on Algebra, Lemma
\ref{more-algebra-lemma-lci-NL} を参照せよ。上の結果を用いると、これがある場合には
局所完全交叉準同型を特徴付けることを示せる。

\begin{lemma}
\label{dpa-lemma-perfect-NL-lci}
$A \to B$ を Noether 環の完全な環準同型とする（More on Algebra, Definition
\ref{more-algebra-definition-pseudo-coherent-perfect}）。次の条件は同値である。
\begin{enumerate}
\item $\NL_{B/A}$ の Tor 振幅は $[-1, 0]$ にある。
\item $\NL_{B/A}$ は $D(B)$ の完全対象で、Tor 振幅は $[-1, 0]$ にある。
\item $A \to B$ は局所完全交叉である（More on Algebra, Definition
\ref{more-algebra-definition-local-complete-intersection}）。
\end{enumerate}
\end{lemma}

\begin{proof}
$B = A[x_1, \ldots, x_n]/I$ と書く。すると $\NL_{B/A}$ は $B$-加群の複体
$$
I/I^2 \longrightarrow \bigoplus B \text{d}x_i
$$
によって表され、$I/I^2$ は次数 $-1$ に置かれる。次数 $0$ の項は有限自由なので、
この複体の Tor 振幅が $[-1, 0]$ にあるための必要十分条件は、$I/I^2$ が平坦 $B$-加群であることである。
More on Algebra, Lemma \ref{more-algebra-lemma-last-one-flat} を参照せよ。
$I/I^2$ は有限 $B$-加群で $B$ は Noether なので、これはさらに $I/I^2$ が
有限局所自由 $B$-加群であることと同値である
（Algebra, Lemma \ref{algebra-lemma-finite-projective}）。
したがって (1) と (2) の同値性は明らかである。また、命題
\ref{dpa-proposition-regular-ideal} を適用し、正則イデアルは Koszul 正則イデアルかつ
擬正則イデアルでもあることに注意すれば、(1) と (3) の同値性も従う。
More on Algebra, Section \ref{more-algebra-section-ideals} を参照せよ。
\end{proof}

\begin{lemma}
\label{dpa-lemma-flat-fp-NL-lci}
$A \to B$ を有限表示の平坦環準同型とする。次の条件は同値である。
\begin{enumerate}
\item $\NL_{B/A}$ の Tor 振幅は $[-1, 0]$ にある。
\item $\NL_{B/A}$ は $D(B)$ の完全対象で、Tor 振幅は $[-1, 0]$ にある。
\item $A \to B$ は syntomic である（Algebra, Definition
\ref{algebra-definition-lci}）。
\item $A \to B$ は局所完全交叉である（More on Algebra, Definition
\ref{more-algebra-definition-local-complete-intersection}）。
\end{enumerate}
\end{lemma}

\begin{proof}
(3) と (4) の同値性は More on Algebra, Lemma
\ref{more-algebra-lemma-syntomic-lci} である。

\medskip\noindent
$A \to B$ が syntomic ならば、余 Cartesian 図式
$$
\xymatrix{
B_0 \ar[r] & B \\
A_0 \ar[r] \ar[u] & A \ar[u]
}
$$
で、$A_0 \to B_0$ が syntomic かつ $A_0$ が Noether であるものが見つかる。
Algebra, Lemmas \ref{algebra-lemma-limit-module-finite-presentation} および
\ref{algebra-lemma-colimit-lci} を参照せよ。補題 \ref{dpa-lemma-perfect-NL-lci} により、
$\NL_{B_0/A_0}$ は Tor 振幅が $[-1, 0]$ にある完全対象である。
More on Algebra, Lemma \ref{more-algebra-lemma-base-change-NL-flat} により、
$\NL_{B/A} = \NL_{B_0/A_0} \otimes_{B_0}^\mathbf{L} B$ についても同じことが成り立つ
（さらに More on Algebra, Lemmas \ref{more-algebra-lemma-pull-tor-amplitude} および
\ref{more-algebra-lemma-pull-perfect} を参照せよ）。これで (3) が (2) を含意すると証明された。

\medskip\noindent
(1) を仮定する。More on Algebra, Lemma
\ref{more-algebra-lemma-base-change-NL-flat} により、体 $k$ へのすべての環準同型
$A \to k$ に対して、$\NL_{B \otimes_A k/k}$ の Tor 振幅は $[-1, 0]$ にある
（More on Algebra, Lemma \ref{more-algebra-lemma-pull-tor-amplitude} を参照せよ）。
したがって補題 \ref{dpa-lemma-perfect-NL-lci} により、$k \to B \otimes_A k$ は
局所完全交叉準同型である。ゆえに定義により $A \to B$ は syntomic である。
これで (1) が (3) を含意すると証明され、証明が完了する。
\end{proof}





\section{Koszul 複体と Tate 分解}
\label{dpa-section-koszul-vs-tate}

\noindent
本節では More on Algebra, Lemma
\ref{more-algebra-lemma-sequence-Koszul-complexes} の結果を、微分次数付き構造と両立する
分割冪を備えた微分次数付き代数の圏に「持ち上げる」
（本節では Koszul 複体を鎖複体として表すが、引用箇所では余鎖複体を用いることに注意せよ）。

\medskip\noindent
$R$ を環とする。$I \subset R$ を $f_1, \ldots, f_r \in R$ によって生成される
イデアルとする。$n \geq 1$ に対して、
$$
K_n = K_{n, \bullet} = R\langle \xi_1, \ldots, \xi_r\rangle
$$
を、$\xi_i$ が次数 $1$ にあり $\text{d}(\xi_i) = f_i^n$ である微分次数付き
Koszul 代数とする。この上には（定義 \ref{dpa-definition-divided-powers-dga} のような）
分割冪構造が一意に存在する。例 \ref{dpa-example-adjoining-odd} を参照せよ。
$m > n$ に対して、遷移写像 $K_m \to K_n$ は $\xi_i$ を
$f_i^{m - n}\xi_i$ に写す、分割冪と両立する微分次数付き代数の写像である。

\begin{lemma}
\label{dpa-lemma-lift-tate-to-koszul}
上の状況で $R$ が Noether 環ならば、各 $n$ に対して $N \geq n$ と写像
$$
K_N \to A \to R/(f_1^N, \ldots, f_r^N)\quad\text{and}\quad A \to K_n
$$
で、次の性質をもつものが存在する。
\begin{enumerate}
\item $(A, \text{d}, \gamma)$ は定義
\ref{dpa-definition-divided-powers-dga} のとおりである。
\item $A \to R/(f_1^N, \ldots, f_r^N)$ は擬同型である。
\item 合成 $K_N \to A \to R/(f_1^N, \ldots, f_r^N)$ は標準写像である。
\item 合成 $K_N \to A \to K_n$ は遷移写像である。
\item $A_0 = R \to R/(f_1^N, \ldots, f_r^N)$ は標準的な全射である。
\item $A$ は $R$ 上の次数付き分割冪多項式代数で、各次数の生成元は有限個である。
\item $A \to K_n$ は分割冪と両立する微分次数付き $R$-代数準同型で、
次数 $0$ のホモロジー上に標準写像
$R/(f_1^N, \ldots, f_r^N) \to R/(f_1^n, \ldots, f_r^n)$。
を誘導する。
\end{enumerate}
条件 (4) は、例 \ref{dpa-example-adjoining-odd} または \ref{dpa-example-adjoining-even} のように、
各次数 $> 0$ で変数 $T$ の有限集合を順次添加することにより、$A$ が $A_0$ から
構成されることを意味する。
\end{lemma}

\begin{proof}
$n$ を固定する。$r = 0$ ならば、単に $A = R$ と選べる。$r > 0$ と仮定する。
More on Algebra, Lemma \ref{more-algebra-lemma-sequence-Koszul-complexes}
（鎖複体の言葉に翻訳したもの）により、
$$
n_{r} > n_{r - 1} > \ldots > n_1 > n_0 = n
$$
を、Koszul 代数の遷移写像 $K_{n_{i + 1}} \to K_{n_i}$（上を参照）が
次数 $> 0$ のホモロジー上に零写像を誘導するように選べる。
$N = n_r$ として補題を証明する。

\medskip\noindent
補題 \ref{dpa-lemma-tate-resolution} の主張と証明における方法と全く同じように $A$ を構成する。
したがって、
$$
A = \colim A(m),\quad\text{and}\quad
A(0) \to A(1) \to A(2) \to \ldots \to R/(f_1^N, \ldots, f_r^N)
$$
となる。これにより性質 (1)、(2)、(5)、(6) が直ちに得られる。
証明を終えるため、$R$-代数写像 $K_N \to A \to K_n$ を構成する。
そのため、次を構成する。
\begin{enumerate}
\item 同型 $A(1) \to K_N = K_{n_r}$,
\item 写像 $A(2) \to K_{n_{r - 1}}$,
\item $\ldots$
\item 写像 $A(r) \to K_{n_1}$,
\item 写像 $A(r + 1) \to K_{n_0} = K_n$, および
\item 写像 $A \to K_n$.
\end{enumerate}
各段階で構成する写像は、両立する分割冪を備えた微分次数付き代数の間の写像である。
各写像は Koszul 代数間の遷移写像を介して前の写像と両立し、次数 $0$ のホモロジー上に
明らかな標準写像を誘導する。

\medskip\noindent
$A(0) = R$ であることを思い出す。$m = 1$ に対して、補題
\ref{dpa-lemma-tate-resolution} の証明では、$T_i$ の次数が $1$ で
$\text{d}(T_i) = f_i^N$ となる
$A(1) = R\langle T_1, \ldots, T_r\rangle$ を選ぶ。
すなわち、$f_i^N$ は $A(0) \to R/(f_1^N, \ldots, f_r^N)$ の核の生成元である。
そこで $A(1) \to K_N = K_{n_r}$ として写像
$$
\varphi_1 : A(1) \longrightarrow K_{n_r},\quad T_i \longmapsto \xi_i
$$
を用いる。これは同型である。

\medskip\noindent
$m = 2$ に対して、補題 \ref{dpa-lemma-tate-resolution} の証明における構成では、
生成元 $e_1, \ldots, e_t \in \Ker(\text{d} : A(1)_1 \to A(1)_0)$ を選ぶ。
次に、$T_i$ の次数が $2$ で $\text{d}(T_i) = e_i$ となる分割冪多項式代数として
$A(2) = A(1)\langle T_1, \ldots, T_t\rangle$
をとる。$\varphi_1(e_i)$ は $K_{n_r}$ のサイクルなので、上での $n_r$ と
$n_{r - 1}$ の選択により、その $K_{n_{r - 1}}$ における像は境界である。
したがって可換図式
$$
\xymatrix{
A(1) \ar[d] \ar[r]_{\varphi_1} & K_{n_r} \ar[d] \\
A(2) \ar[r]^{\varphi_2} & K_{n_{r - 1}}
}
$$
を、$T_i$ を、その境界が $\varphi_1(e_i)$ の像となる次数 $2$ の元に写すことで構成できる。
分割冪多項式代数の普遍性により写像 $\varphi_2$ は存在し、微分および分割冪と両立する。

\medskip\noindent
$m = 2, \ldots, r$ に対して、代数 $A(m)$ と写像
$\varphi_m : A(m) \to K_{n_{r + 1 - m}}$ を全く同じ方法で構成する。

\medskip\noindent
写像 $A(r) \to K_{n_1}$ が与えられたとき、合成
$H_r(A(r)) \to H_r(K_{n_1}) \to H_r(K_{n_0}) \subset (K_{n_0})_r$
は零である。したがって、$A(r + 1)$ の新しい多項式生成元を零に写すことで、
これを $A(r + 1) \to K_{n_0} = K_n$ に延長できる。

\medskip\noindent
$A(r + 1) \to K_{n_0} = K_n$ を構成した後、新しい多項式生成元を零に写すことにより、
唯一可能な仕方で $A(r + 2), A(r + 3), \ldots$ に延長できる。
これで証明が完了する。
\end{proof}

\begin{remark}
\label{dpa-remark-pro-system-koszul}
上の状況で $R$ が Noether 環ならば、整数の列
$1 = n_0 < n_1 < n_2 < \ldots $ を帰納的に選んで、
$i = 1, 2, 3, \ldots$ に対して補題
\ref{dpa-lemma-lift-tate-to-koszul} のような写像
$K_{n_i} \to A_i \to R/(f_1^{n_i}, \ldots, f_r^{n_i})$
および $A_i \to K_{n_{i - 1}}$ が存在するようにできる。
合成 $A_{i + 1} \to K_{n_i} \to A_i$ を $A_{i + 1} \to A_i$ と表す。
すると図式
$$
\xymatrix{
K_{n_1} \ar[d] &
K_{n_2} \ar[d] \ar[l] &
K_{n_3} \ar[d] \ar[l] &
\ldots \ar[l] \\
A_1 \ar[d] &
A_2 \ar[l] \ar[d] &
A_3 \ar[l] \ar[d] &
\ldots \ar[l] \\
K_1 &
K_{n_1} \ar[l] &
K_{n_2} \ar[l] &
\ldots \ar[l]
}
$$
は可換である。このようにして、両立する分割冪を備えた微分次数付き
$R$-代数の圏において、逆系 $(K_n)$ と $(A_n)$ は pro-同型であることが分かる。
\end{remark}

\begin{lemma}
\label{dpa-lemma-compute-cohomology-adjoin-variable}
$(A, \text{d}, \gamma)$、$d \geq 1$、$f \in A_{d - 1}$ および
$A\langle T \rangle$ を補題 \ref{dpa-lemma-extend-differential} と同様とする。
\begin{enumerate}
\item $d = 1$ ならば、長完全列
$$
\ldots \to H_0(A) \xrightarrow{f} H_0(A) \to H_0(A\langle T \rangle) \to 0
$$
が存在する。
\item $d = 2$ ならば、有界なスペクトル系列
$(E_1)_{i, j} = H_{j - i}(A) \cdot T^{[i]}$
が存在し、$H_{i + j}(A\langle T \rangle)$ に収束する。その微分
$(d_1)_{i, j} : H_{j - i}(A) \cdot T^{[i]} \to
H_{j - i + 1}(A) \cdot T^{[i - 1]}$
は $\xi \cdot T^{[i]}$ を $f \xi \cdot T^{[i - 1]}$ の類へ送る。
\item 必要に応じて、他の次数についてここにさらに追加する。
\end{enumerate}
\end{lemma}

\begin{proof}
$d = 1$ ならば、複体の短完全列
$$
0 \to A \to A\langle T \rangle \to A \cdot T \to 0
$$
があり、これから (1) が直ちに従う。$d = 2$ ならば、
$A\langle T \rangle$ を部分複体
$$
F^pA\langle T \rangle = \bigoplus\nolimits_{i \leq p} A \cdot T^{[i]}
$$
をもつフィルター付き鎖複体とみなす。
Homology, Section \ref{homology-section-filtered-complex} のスペクトル系列を
（鎖複体に読み替えて）適用すると (2) を得る。
\end{proof}

\noindent
次の補題は後で用いる。

\begin{lemma}
\label{dpa-lemma-construct-some-maps}
上の状況で、任意の $n \geq t \geq 1$ に対し、ある $N > n$ と写像
$$
K_t \longrightarrow K_n \otimes_R K_t
$$
が左微分次数付き $K_N$-加群の導来圏に存在し、これを乗法写像と
いずれの向きに合成しても推移写像になる。
\end{lemma}

\begin{proof}
まず $r = 1$ の場合を証明する。$f = f_1$ とおく。
$K_t = R\langle x \rangle$、
$K_n = R\langle y \rangle$、$K_N = R\langle z \rangle$ と書く。
ここで $x,y,z$ は次数 $1$ であり、
$\text{d}(x) = f^t$、$\text{d}(y) = f^n$、$\text{d}(z) = f^N$ である。
任意の $N > t$ に対して擬同型
$$
B_{N, t} = R\langle x, z, u \rangle
\longrightarrow
K_t,\quad
x \mapsto x,\quad
z \mapsto f^{N - t}x,\quad
u \mapsto 0
$$
が存在すると主張する。ここで左辺は、変数 $x,z$ の次数が $1$、
変数 $u$ の次数が $2$ である分割冪多項式代数を表し、
$\text{d}(x) = f^t$、$\text{d}(z) = f^N$、
$\text{d}(u) = z - f^{N - t}x$ とする。この主張を証明するため、
$H_*(R\langle x, z\rangle)$ の次の三つの部分加群が一致することに注意する。
\begin{enumerate}
\item $H_*(R\langle x, z\rangle) \to H_*(K_t)$ の核、
\item 写像
$z - f^{N - t}x : H_*(R\langle x, z\rangle) \to H_*(R\langle x, z\rangle)$
の像、および
\item 写像
$z - f^{N - t}x : H_*(R\langle x, z\rangle) \to H_*(R\langle x, z\rangle)$
の核。
\end{enumerate}
この事実は直接計算で証明できる\footnote{ヒント：
$z' = z - f^{N - t}x$ とおくと、
$R\langle x, z\rangle = R\langle x, z'\rangle$ かつ
$\text{d}(z') = 0$ である。さらに、写像
$R\langle x, z'\rangle \to K_t$ は $z'$ を零にする写像である。}が、
計算は省略する。次に補題
\ref{dpa-lemma-compute-cohomology-adjoin-variable} の (2) を適用すれば、
この主張が従う。

\medskip\noindent
$z$ を $z$ へ送る微分次数付き $R$-代数の準同型
$K_N \to B_{N, t}$ によって、$B_{N, t} \to K_t$ を左微分次数付き
$K_N$-加群の擬同型とみなせる。補題の主張にある矢印を定義するため、準同型
$$
B_{N, t} = R\langle x, z, u \rangle \to K_n \otimes_R K_t,\quad
x \mapsto 1 \otimes x,\quad
z \mapsto f^{N - n}y \otimes 1,\quad
u \mapsto - f^{N - n - t}y \otimes x
$$
を用いる。これは $N \geq n + t$ と仮定すれば意味をもつ。
乗法写像との前後いずれの合成も推移写像になることは、簡単な計算で確かめられる。

\medskip\noindent
$r > 1$ の場合は、各 Koszul 代数を一変数 Koszul 代数のテンソル積として書き、
上の構成を適用する。すなわち、
$$
K_t = R\langle x_1, \ldots, x_r\rangle =
R\langle x_1\rangle \otimes_R \ldots \otimes_R R\langle x_r\rangle
$$
と書く。ここで $x_i$ は次数 $1$ で、$\text{d}(x_i) = f_i^t$ である。
$r > 1$ の場合には
$$
B_{N, t} =
R\langle x_1, z_1, u_1 \rangle
\otimes_R \ldots \otimes_R
R\langle x_r, z_r, u_r \rangle
$$
を用いる。ここで $x_i,z_i$ は次数 $1$、$u_i$ は次数 $2$ であり、
$\text{d}(x_i) = f_i^t$、$\text{d}(z_i) = f_i^N$、
$\text{d}(u_i) = z_i - f_i^{N - t}x_i$ である。テンソル積写像
$B_{N, t} \to K_t$ は、自由 $R$-加群の上に有界な複体間の擬同型の
テンソル積なので、擬同型である。最後に、写像
$$
B_{N, t}
\to
K_n \otimes_R K_t =
R\langle y_1, \ldots, y_r\rangle \otimes_R
R\langle x_1, \ldots, x_r\rangle
$$
を $r = 1$ の場合に構成した写像のテンソル積として定義する。
あるいは、規則 $x_i \mapsto 1 \otimes x_i$、
$z_i \mapsto f_i^{N - n}y_i \otimes 1$、
$u_i \mapsto - f_i^{N - n - t}y_i \otimes x_i$ によって直接定義してもよい。
これらは $N \geq n + t$ ならば意味をもつ。詳細は省略する。
\end{proof}
